\documentclass[11pt,letterpaper]{article}

% --- PAQUETES DE DISEÑO ---
\usepackage[utf8]{inputenc}
% es-nodecimaldot evita el conflicto de babel-español con siunitx
\usepackage[spanish, es-tabla, es-nodecimaldot]{babel}
\usepackage[T1]{fontenc}
\usepackage{lmodern}
\usepackage{microtype}
\usepackage{geometry}
\usepackage{graphicx}
\usepackage{float}
\usepackage{url}
\usepackage{hyperref}
\usepackage{xurl}          % corta las URLs largas para que no se salgan del margen
\usepackage{listings}
\usepackage{amsmath}
\usepackage{amssymb}
\usepackage{booktabs}
\usepackage{siunitx}
\usepackage{subcaption}
\usepackage[authoryear,round]{natbib}
\usepackage{tikz}
\usepackage{pgfplots}
\usetikzlibrary{shapes.geometric, arrows.meta, positioning, calc}
\usepackage{adjustbox}
\geometry{top=2cm, bottom=2.5cm, left=2.2cm, right=2.2cm}
% Evita que identificadores largos en \texttt se salgan del margen.
\emergencystretch=2.5em

% --- COLORES ---
\usepackage{xcolor}
\definecolor{mainBlue}{RGB}{0, 32, 96}
\definecolor{lightGray}{RGB}{240, 240, 240}

% --- TÍTULOS Y SECCIONES ---
\usepackage{titlesec}
\titleformat{\section}{\large\bfseries\color{mainBlue}}{\thesection.}{0.5em}{}[\titlerule]

% --- CAJAS Y CÓDIGO ---
\usepackage{tcolorbox}

\lstset{
    basicstyle=\ttfamily\footnotesize,
    backgroundcolor=\color{lightGray!50},
    frame=single,
    rulecolor=\color{mainBlue!20},
    columns=fullflexible,
    breaklines=true,
    breakatwhitespace=true,
    showstringspaces=false,
    xleftmargin=10pt,
    xrightmargin=6pt,
    inputencoding=utf8,
}

\newtcolorbox{analisisBox}[1]{
    colback=lightGray!30, colframe=mainBlue, fonttitle=\bfseries,
    title=#1, arc=0pt, outer arc=0pt, width=\linewidth,
    left=4pt, right=4pt
}


% siunitx: separador de miles con espacio fino, coma decimal
\sisetup{group-separator={\,}, output-decimal-marker={,}, group-minimum-digits=4, group-digits=integer}

\hypersetup{colorlinks=true, linkcolor=mainBlue, urlcolor=mainBlue, citecolor=mainBlue,
            breaklinks=true}

% --- INICIO DEL DOCUMENTO ---
\begin{document}
% babel-spanish activa "<" y ">" como atajos para comillas latinas;
% eso choca con listings y con la sintaxis >={...} de TikZ.
\shorthandoff{<>"}

% --- ENCABEZADO COMPACTO ---
\begin{center}
    {\color{mainBlue}\hrule height 1.5pt}
    \vspace{0.4cm}
    {\Large \textbf{\texttt{Ejercicio No.\ 3:}\\
    \vspace{0.4cm}
    Modelo QGeoidCOL2023\\
    \vspace{0.15cm}
    \large Interpolación de la anomalía de altura y determinación de la altura normal}} \\
    \vspace{0.2cm}
    Grupo 2 \\
    \vspace{0.4cm}
    {\large \textit{Sistemas de Posicionamiento Geodésico}}\\
    {\large \textit{Maestría en Ciencias de la Información y las Comunicaciones \\ Énfasis en Geomática\\ 2026-3}} \\
    \vspace{0.4cm}
    {\color{mainBlue!30}\hrule height 0.8pt}
    \vspace{0.2cm}

    \begin{tabular}{cc}
        \textbf{Autor 1} & \textbf{Autor 2} \\[0.2cm]
        Angie Tatiana Díaz & Natalia Lizeth Rodríguez \\[0.1cm]
        \footnotesize \texttt{atdiazl@udistrital.edu.co} &
        \footnotesize \texttt{nlrodriguezv@udistrital.edu.co} \\[0.5cm]
        \textbf{Autor 3} & \textbf{Autor 4} \\[0.2cm]
        Julián Santiago López & Argenis Alexandra Daza Roa \\[0.1cm]
        \footnotesize \texttt{jslopezc@udistrital.edu.co} &
        \footnotesize \texttt{aadazar@udistrital.edu.co}
    \end{tabular}

    \vspace{0.5cm}
    \small \textbf{Docente:} John Alexander Orduña Díaz \quad|\quad \textbf{Fecha:} \today
    \vspace{0.2cm}
    {\color{mainBlue}\hrule height 1.5pt}
\end{center}

\begin{abstract}
\noindent
Se desarrolló en MATLAB una herramienta con interfaz gráfica
(\texttt{appQGeoidCOL2023.m}) que lee y valida la malla oficial \textbf{QGeoidCOL2023},
interpola la anomalía de altura $\zeta$ en cualquier punto dentro de su cobertura y
calcula la altura normal $H^N=h-\zeta$. La malla (\num{32224} nodos, $212\times152$, paso
de $5'$, $\approx\SI{9.2}{km}$) superó todos los controles de calidad: datos numéricos y
finitos, sin duplicados ni nodos faltantes, rectangular y regular con desviaciones del
paso del orden de $10^{-14}$ grados. Se implementaron y compararon \textbf{seis métodos}:
vecino más cercano, lineal por triángulos, bilineal, bicúbica (convolución cúbica),
IDW y kriging ordinario. Cada uno reporta de forma explícita los nodos que usa, sus pesos y sus
contribuciones, y verifica que los pesos sumen uno. Ninguno extrapola fuera de la malla.
Las implementaciones del vecino más cercano, la bilineal y la bicúbica coinciden con
\texttt{interp2} de MATLAB a nivel de $10^{-14}$\,m. En cinco puntos de control con altura
elipsoidal, altura nivelada y valores de GeoCol2004 y EGM2008 (archivo de puntos del grupo),
los métodos (excepto el vecino más cercano) difieren entre sí entre \SI{21}{mm} y \SI{83}{mm}; el vecino más
cercano se aparta hasta \SI{25}{cm}. Una validación cruzada por decimación de la malla ubica a la
bicúbica como el método de menor error (RMSE de \SI{58}{mm} con paso de prueba de $10'$),
seguida de la bilineal y la lineal por triángulos (\SI{87}{mm} y \SI{89}{mm}). La diferencia
QGeoidCOL2023 $-$ GeoCol2004 ($\zeta-N$) es negativa en los cinco puntos, de \SI{-0.73}{m} a
\SI{-4.07}{m} (media \SI{-2.14}{m}), y no se explica por la separación geoide--cuasigeoide sino por
el sesgo de GeoCol2004, cuya diferencia media con la nivelación en estos puntos (\SI{-2.1}{m}) es del
orden de la documentada para ese modelo. Frente a EGM2008 (ondulación geoidal), la diferencia va de
\SI{-21}{cm} a \SI{+70}{cm} (RMS de \SI{36}{cm}). Como control altimétrico, en los tres puntos
bajos (P1, P2 y P5) la altura normal calculada con QGeoidCOL2023 difiere de la altura nivelada en
\SI{-1.8}{cm}, \SI{+15.9}{cm} y \SI{+1.0}{cm} (RMS de \SI{9}{cm}), frente a \SI{47}{cm} con EGM2008 y
\SI{112}{cm} con GeoCol2004. En los dos puntos de montaña (P3 y P4) las diferencias con la
nivelación son de \SI{+65}{cm} y \SI{-62}{cm}, de magnitud similar a las de EGM2008.
\end{abstract}

\tableofcontents
\vspace{0.6cm}

% =====================================================================
\section{Introducción}
% =====================================================================

El posicionamiento GNSS entrega alturas \emph{elipsoidales} $h$, referidas a una
superficie matemática (el elipsoide GRS80 en MAGNA-SIRGAS) que no tiene significado
físico: el agua no fluye necesariamente de una $h$ mayor a una menor. Las alturas de uso
práctico en ingeniería, hidrología y cartografía deben estar referidas al campo de
gravedad terrestre. Para convertir una $h$ en una altura física se necesita conocer la
separación entre el elipsoide y la superficie de referencia vertical correspondiente: la
\emph{ondulación geoidal} $N$ si la altura es ortométrica, o la \emph{anomalía de altura}
$\zeta$ si la altura es normal \citep{hofmannmoritz2006,torgemuller2012}.

\textbf{QGeoidCOL2023} es el modelo cuasigeoidal oficial de Colombia, adoptado por el
IGAC en 2026 como reemplazo de GeoCol2004 \citep{igacnoticia}. Es un cuasigeoide gravimétrico
calculado por el DGFI-TUM con el IGAC a partir de gravedad terrestre y aerotransportada, altimetría
satelital sobre el mar (DTU21GRA), un modelo global de fondo y modelos topográficos de muy alta
resolución, con funciones de base radial esféricas \citep{liu2024,isg}. Sus anomalías de altura se
calculan en el sistema de marea cero, respecto al elipsoide GRS80 y con el potencial convencional
$W_0=\SI{62636853.4}{m^2/s^2}$ del Sistema Internacional de Referencia de Alturas (IHRS), y se
distribuyen en una malla de $5'$ \citep{liu2024,isg}. Con 3025 puntos GPS/nivelación, los autores
reportan una desviación estándar de las diferencias de \SI{21.3}{cm} y de \SI{15.8}{cm} tras ajustar
una superficie de corrección de siete parámetros \citep{liu2024}. Con el modelo, la altura normal de
un punto se obtiene directamente de la altura elipsoidal medida con GNSS:
\begin{equation}
H^N = h - \zeta.
\label{eq:hn}
\end{equation}
Este procedimiento, conocido como \emph{nivelación GNSS}, reemplaza en muchos casos la
nivelación geométrica, que es más costosa y lenta. Su exactitud depende de la calidad del
modelo y del método con que se interpola $\zeta$ entre los nodos de la malla: con un paso
de $5'$ ($\approx\SI{9}{km}$), en zonas de topografía abrupta la forma de interpolar puede
cambiar el resultado en centímetros.

En este ejercicio se implementó una herramienta que carga y valida la malla, interpola
$\zeta$ con seis métodos, reporta los nodos, pesos y contribuciones de cada uno, calcula
$H^N$, controla la cobertura para no extrapolar y compara los métodos entre sí. También se
contrastan los resultados con los modelos GeoCol2004 y EGM2008 en cinco puntos del
territorio colombiano.

% =====================================================================
\section{Objetivos}
% =====================================================================
\subsection{Objetivo general}
Implementar una herramienta computacional para la lectura, validación, interpolación y
representación gráfica del modelo QGeoidCOL2023, que permita analizar la anomalía de
altura $\zeta$ y su aplicación en la relación entre la altura elipsoidal $h$ y la altura
normal $H^N$.

\subsection{Objetivos específicos}
\begin{itemize}
    \item Cargar el archivo \texttt{QGeoidCOL2023.TXT}, reconstruir la malla y validar su
    estructura: datos finitos, coordenadas válidas, duplicados, nodos faltantes,
    regularidad, dimensiones, resolución y cobertura.
    \item Implementar seis métodos de interpolación (vecino más cercano, lineal por
    triángulos, bilineal, bicúbica, IDW y kriging ordinario) que reporten de forma explícita
    los nodos usados, sus pesos y sus contribuciones, y verificar que la suma de pesos sea uno.
    \item Impedir la extrapolación fuera de la cobertura y declarar, sin sustituirlo por
    otro, cuándo un método no puede aplicarse por falta de soporte espacial.
    \item Calcular $H^N=h-\zeta$ conservando la precisión numérica y declarando el criterio
    de redondeo de la presentación.
    \item Construir una visualización general de la malla y una vista local del campo
    interpolado, con los nodos usados y la geometría de soporte de cada método.
    \item Cuantificar las diferencias entre métodos en metros, centímetros y milímetros en
    cinco puntos distribuidos en Colombia, y evaluar su desempeño mediante validación
    cruzada y un análisis de sensibilidad de parámetros.
    \item Comparar QGeoidCOL2023 con GeoCol2004 y EGM2008 en los mismos cinco puntos, tras
    verificar la compatibilidad conceptual de las magnitudes de cada modelo.
\end{itemize}

% =====================================================================
\section{Datos y control de calidad}
\label{sec:datos}
% =====================================================================

\subsection{Estructura del archivo}
El archivo \texttt{QGeoidCOL2023.TXT} no tiene encabezado y contiene \num{32224} líneas,
cada una con tres columnas numéricas separadas por espacios (Tabla~\ref{tab:notacion}).
Viene ordenado por filas de latitud decreciente (de norte a sur) y, dentro de cada fila,
por longitud creciente (de oeste a este). La herramienta no depende de ese orden, porque
reconstruye la malla a partir de los índices de cada nodo.

\begin{table}[H]
\centering
\caption{Columnas de QGeoidCOL2023.TXT y notación adoptada.}
\label{tab:notacion}
\begin{tabular}{@{}clll@{}}
\toprule
Col. & Variable & Notación & Unidad \\
\midrule
1 & Latitud geodésica del nodo & $\varphi$ & grados decimales (+N) \\
2 & Longitud geodésica del nodo & $\lambda$ & grados decimales ($-$ al oeste) \\
3 & Anomalía de altura & $\zeta$ & m \\
\bottomrule
\end{tabular}
\end{table}

\subsection{Controles de calidad}
La función \texttt{cargarYValidarMalla} aplica los controles en el orden de la
Tabla~\ref{tab:qc}. Cada control remueve los registros que detecta antes de pasar al
siguiente y deja registrado el criterio que usó, de modo que un conteo en cero significa
``se verificó y no había nada que corregir'' y no ``no se revisó''.

\begin{table}[H]
\centering
\caption{Controles de calidad aplicados sobre QGeoidCOL2023.TXT.}
\label{tab:qc}
\renewcommand{\arraystretch}{1.25}
\begin{tabular}{@{}p{3.4cm}p{7.6cm}r@{}}
\toprule
Control & Criterio & Resultado \\
\midrule
Datos numéricos & Cada línea tiene exactamente 3 números (\texttt{sscanf}) & 0 líneas inválidas \\
Valores finitos & Ningún \texttt{NaN} ni \texttt{Inf} & 0 \\
Rango geográfico & $|\varphi|\le\ang{90}$ y $|\lambda|\le\ang{180}$ & 0 fuera de rango \\
Duplicados & Misma posición de nodo repetida (y, aparte, con $\zeta$ distinta) & 0 (0) \\
Regularidad & Paso constante en $\varphi$ y en $\lambda$ (tolerancia $10^{-7}$ grados) & Sí \\
Dimensiones & Valores únicos de $\varphi$ y $\lambda$ & $212\times152=\num{32224}$ \\
Nodos faltantes & Posiciones de la malla rectangular sin dato & 0 \\
Plausibilidad & $|\zeta|>\SI{150}{m}$ (no ocurre en la Tierra) & 0 \\
\midrule
\multicolumn{2}{l}{\textbf{Nodos válidos}} & \textbf{\num{32224} / \num{32224}} \\
\bottomrule
\end{tabular}
\end{table}

Las coordenadas del archivo se escriben con representación binaria imprecisa (por ejemplo,
\texttt{12.800000000000001}). Para agruparlas se usó \texttt{uniquetol} con tolerancia de
$10^{-9}$ grados ($\approx\SI{0.1}{mm}$), que conserva los valores reales del archivo en
lugar de redondearlos de antemano. Así, las desviaciones que se reportan corresponden a los
datos y no a un artificio del redondeo: la desviación máxima del paso es de
$\num{2.2e-15}$ grados en latitud y $\num{2.8e-14}$ grados en longitud, y la de las
coordenadas frente a la malla nominal es de $\num{1.4e-14}$ grados.

\subsection{Resolución, cobertura y estadísticas}
\begin{table}[H]
\centering
\caption{Características de la malla reconstruida.}
\label{tab:malla}
\begin{tabular}{@{}ll@{}}
\toprule
Característica & Valor \\
\midrule
Dimensiones & 212 filas (latitud) $\times$ 152 columnas (longitud) \\
Paso & $\Delta\varphi=\Delta\lambda=\ang{0.0833333333}=5'$ \\
Resolución métrica (latitud media, GRS80) & \SI{9.215}{km} (N--S) $\times$ \SI{9.254}{km} (E--O) \\
Cobertura en latitud & $\ang{-4.783333}$ a $\ang{12.800000}$ \\
Cobertura en longitud & $\ang{-79.300000}$ a $\ang{-66.716667}$ \\
Extensión aproximada & \SI{1944}{km} (N--S) $\times$ \SI{1397}{km} (E--O) \\
$\zeta$ mínima & \SI{-39.949}{m} en $(\ang{12.7167},\ang{-66.7167})$ \\
$\zeta$ máxima & \SI{33.425}{m} en $(\ang{3.8000},\ang{-75.9667})$ \\
$\zeta$ media y desviación estándar & \SI{6.927}{m} y \SI{13.685}{m} \\
\bottomrule
\end{tabular}
\end{table}

La Figura~\ref{fig:malla} muestra el campo de $\zeta$. Los valores máximos (más de
\SI{30}{m}) siguen el eje de las cordilleras Central y Occidental, mientras que los
mínimos (hasta \SI{-40}{m}) aparecen en el Caribe nororiental. Esto es coherente con
el hecho de que las anomalías de altura reflejan la distribución de masas. El histograma
es bimodal: una moda cerca de \SI{20}{m}, que corresponde a la Orinoquía, la Amazonía y la
región andina, y otra cerca de \SI{-5}{m}, que corresponde a la región Caribe.

\begin{figure}[H]
\centering
\begin{subfigure}[t]{0.51\linewidth}
\centering
\includegraphics[width=\linewidth]{figuras/malla_general.png}
\caption{Campo de $\zeta$ con isolíneas cada \SI{5}{m}.}
\end{subfigure}
\hfill
\begin{subfigure}[t]{0.45\linewidth}
\centering
\includegraphics[width=\linewidth]{figuras/histograma_zeta.png}
\caption{Histograma de los \num{32224} nodos.}
\end{subfigure}
\caption{Malla QGeoidCOL2023 reconstruida. El rectángulo negro es el límite de cobertura:
fuera de él la herramienta no interpola.}
\label{fig:malla}
\end{figure}

\subsection{Puntos de control}
Para las comparaciones se usaron cinco puntos del archivo \texttt{5 PUNTOS CALCULADOS COLOMBIA.xlsx}
entregado al grupo (Tabla~\ref{tab:puntos}). Para cada punto el archivo trae la altura elipsoidal $h$,
la altura $H$ de la nivelación geométrica (altura nivelada y ajustada por mínimos cuadrados, con el
año de cálculo; no se clasifica como ortométrica ni como normal) y los valores de cada modelo: la
ondulación geoidal $N$ de GeoCol2004 y de EGM2008 y la anomalía de altura $\zeta$ de QGeoidCOL2023.
Según el grupo, los valores de GeoCol2004 y de QGeoidCOL2023 se extrajeron del software MAGNA 7.0 y los de
EGM2008 de la calculadora de UNAVCO/EarthScope. Los puntos están en la Guajira (P1), Córdoba (P2),
Norte de Santander, en la cordillera Oriental (P3), Bogotá (P4) y los llanos de Arauca, en la Orinoquía
(P5), según su ubicación aproximada (Figura~\ref{fig:puntos}). Cubren alturas elipsoidales de
\SI{35}{m} a \SI{3229}{m} y valores de $\zeta$ de \SI{-13.3}{m} a \SI{+23.4}{m}, frente a un rango de
\SIrange{-39.9}{33.4}{m} de la malla. Quedan, sin embargo, todos al norte de \ang{4.6} N: no hay puntos
en la Amazonía, el Pacífico ni el sur del país, de modo que las conclusiones no deben extrapolarse
a todo el territorio (Sección~\ref{sec:analisis}).

\begin{table}[H]
\centering
\caption{Puntos de control (archivo \texttt{5 PUNTOS CALCULADOS COLOMBIA.xlsx}).}
\label{tab:puntos}
\small
\begin{tabular}{@{}llrrrc@{}}
\toprule
Id & Nomenclatura & $\varphi$ (grados) & $\lambda$ (grados) & $h$ (m) & $H$ nivelada (m) \\
\midrule
P1 & 32-TN-7 & \num{11.4519608895} & \num{-72.5744326992} & \num{35.134} & \num{48.43433} (2021) \\
P2 & 23555011 & \num{8.4311028480} & \num{-75.5530108810} & \num{98.266} & \num{95.57344} (2021) \\
P3 & GPS-S-T-144 & \num{7.3110555110} & \num{-72.8720174438} & \num{3228.966} & \num{3216.550} (2005) \\
P4 & (EA)-MACARENA-CD-551-1 & \num{4.6149699549} & \num{-74.0628738019} & \num{2725.381} & \num{2702.640} (2001) \\
P5 & 81794007 & \num{6.2760436758} & \num{-71.6514462724} & \num{185.792} & \num{184.956} (2019) \\
\bottomrule
\end{tabular}

\vspace{0.1cm}
\footnotesize Entre paréntesis, el año de cálculo de la nivelación que figura en el archivo.
Todos los puntos están dentro de la cobertura de la malla y ninguno cae en una celda de borde.
\end{table}

\noindent\textbf{Marco de referencia y fuentes.} Las coordenadas y las alturas elipsoidales $h$ están
referidas a MAGNA-SIRGAS, y la época de $h$ es la misma de las coordenadas. Los valores de GeoCol2004 y de
QGeoidCOL2023 se extrajeron de MAGNA 7.0 (versión 7.0, la primera del software) y los de EGM2008 de la calculadora de
UNAVCO. Los valores del archivo se tratan como dato de entrada del grupo y se verifican de forma
independiente en la Sección~\ref{sec:modelos}.

\begin{figure}[H]
\centering
\includegraphics[width=0.62\linewidth]{figuras/malla_puntos.png}
\caption{Ubicación de los cinco puntos de control sobre la malla QGeoidCOL2023.}
\label{fig:puntos}
\end{figure}

% =====================================================================
\section{Marco conceptual}
\label{sec:marco}
% =====================================================================

\subsection{Superficies de referencia y alturas}
\label{sec:alturas}
\begin{itemize}
    \item \textbf{Elipsoide} (GRS80 en MAGNA-SIRGAS): superficie geométrica de referencia
    para las coordenadas $(\varphi,\lambda,h)$. La altura elipsoidal $h$ es la distancia a
    lo largo de la normal al elipsoide y es lo que mide directamente el GNSS.
    \item \textbf{Geoide}: superficie equipotencial del campo de gravedad terrestre que mejor se
    ajusta al nivel medio del mar. La altura \textbf{ortométrica} $H$ se mide desde el
    geoide a lo largo de la línea de la plomada, y su separación del elipsoide es la
    ondulación geoidal $N$. Para calcular $H$ se necesita la gravedad media a lo largo de la plomada,
    \emph{dentro} de las masas topográficas, que no puede observarse y debe suponerse
    (por ejemplo, con la aproximación de Helmert).
    \item \textbf{Cuasigeoide}: superficie que no es equipotencial y que está a una
    distancia $\zeta$ del elipsoide. La altura \textbf{normal} $H^N$ de Molodensky se calcula
    con la gravedad \emph{normal} media, que se obtiene de la fórmula de la gravedad normal sin
    ninguna hipótesis sobre la densidad de la corteza. En este punto está su principal ventaja
    frente a la altura ortométrica \citep{hofmannmoritz2006}.
\end{itemize}

Las dos descomposiciones de la altura elipsoidal son
\begin{equation}
h = H + N = H^N + \zeta,
\label{eq:alturas}
\end{equation}
de donde la separación entre geoide y cuasigeoide resulta
\begin{equation}
N-\zeta = H^N - H \approx \frac{\Delta g_B}{\bar{\gamma}}\,H,
\label{eq:nzeta}
\end{equation}
con $\Delta g_B$ la anomalía de Bouguer y $\bar\gamma$ la gravedad normal media
\citep{hofmannmoritz2006}. En zonas bajas esta separación es de milímetros, pero en
montañas con anomalías de Bouguer de cientos de mGal puede alcanzar decímetros. Por eso
\textbf{$N$ y $\zeta$ no son magnitudes intercambiables} y compararlas directamente
(Sección~\ref{sec:modelos}) exige declarar esa diferencia conceptual.

\subsection{Interpolación sobre la malla: notación común}
Para un punto $P=(\varphi,\lambda)$, todos los métodos implementados producen una
combinación lineal de los valores de la malla:
\begin{equation}
\hat\zeta(P)=\sum_{i\in S(P)} w_i\,\zeta_i,
\qquad \sum_{i\in S(P)} w_i = 1,
\label{eq:combinacion}
\end{equation}
donde $S(P)$ es el conjunto de nodos de soporte y $w_i$ sus pesos. La
\emph{contribución} de cada nodo es $w_i\zeta_i$. Que los pesos sumen uno garantiza que
una superficie constante se reproduzca exactamente, y la herramienta lo verifica en cada
cálculo. Para los métodos sobre la malla se define la celda
$[\varphi_i,\varphi_{i+1}]\times[\lambda_j,\lambda_{j+1}]$ que contiene a $P$ y sus
coordenadas normalizadas
\begin{equation}
u=\frac{\lambda-\lambda_j}{\lambda_{j+1}-\lambda_j},\qquad
v=\frac{\varphi-\varphi_i}{\varphi_{i+1}-\varphi_i},\qquad u,v\in[0,1].
\label{eq:uv}
\end{equation}
Las esquinas de la celda se denotan SW $(i,j)$, SE $(i,j+1)$, NE $(i+1,j+1)$ y NW
$(i+1,j)$.

\subsection{Vecino más cercano}
Asigna a $P$ el valor del nodo más próximo: $|S|=1$ y $w=1$. La distancia se mide en
metros sobre la proyección Origen-Nacional, aunque en una malla regular el nodo más cercano
siempre es una de las cuatro esquinas de la celda, así que el resultado coincide con
\texttt{interp2(...,'nearest')}. El campo resultante es discontinuo, con escalones en
las mediatrices entre nodos.

\subsection{Lineal por triángulos}
La celda se divide en dos triángulos y, dentro del triángulo que contiene a $P$, $\zeta$
se interpola sobre un plano. Los pesos son las coordenadas baricéntricas:
\begin{equation}
(w_{SW},w_{SE},w_{NE}) = (1-u,\;u-v,\;v)\quad\text{si } u\ge v;\qquad
(w_{SW},w_{NE},w_{NW}) = (1-v,\;u,\;v-u)\quad\text{si } u<v.
\label{eq:triangulos}
\end{equation}
\textbf{Decisión metodológica:} la celda se divide \emph{siempre} por la diagonal SW--NE.
En una malla cuadrada los cuatro vértices de cada celda son cocirculares, por lo que la
triangulación de Delaunay no es única: \texttt{scatteredInterpolant} elegiría la diagonal
según detalles numéricos y el resultado dejaría de ser reproducible. El campo es continuo
($C^0$), pero tiene quiebres sobre las aristas.

\subsection{Bilineal}
Usa las cuatro esquinas de la celda:
\begin{equation}
\hat\zeta = (1-u)(1-v)\,\zeta_{SW} + u(1-v)\,\zeta_{SE} + uv\,\zeta_{NE} + (1-u)v\,\zeta_{NW}.
\label{eq:bilineal}
\end{equation}
Es el método habitual para consultar mallas geoidales oficiales. Es continuo, exacto en
los nodos y fácil de verificar a mano, pero la derivada es discontinua en las aristas
de las celdas.

\subsection{Bicúbica (convolución cúbica)}
Usa los $4\times4=16$ nodos alrededor de la celda, con pesos separables
$w_{kl}=w^{\varphi}_k\,w^{\lambda}_l$ obtenidos del núcleo de convolución cúbica de
\citet{keys1981}:
\begin{equation}
W(s)=\begin{cases}
(a+2)|s|^3-(a+3)|s|^2+1, & |s|\le1,\\
a|s|^3-5a|s|^2+8a|s|-4a, & 1<|s|<2,\\
0, & |s|\ge2,
\end{cases}
\qquad a=-0{,}5,
\label{eq:keys}
\end{equation}
evaluado en las distancias normalizadas $s=\{u+1,\,u,\,1-u,\,2-u\}$ (y análogamente con
$v$). Con $a=-0{,}5$ el núcleo es el mismo que usa \texttt{interp2(...,'cubic')} en mallas
uniformes, lo que permitió validar la implementación (Sección~\ref{sec:autoverif}). Algunos
pesos son negativos, por los lóbulos del núcleo, pero cada vector unidimensional suma uno.
\textbf{Limitación:} necesita un nodo adicional a cada lado de la celda. En las celdas del
borde exterior de la malla el método se declara \emph{no aplicable}, con su causa, y no se
reemplaza por otro método.

\subsection{IDW (distancia inversa ponderada)}
\begin{equation}
w_i=\frac{d_i^{-p}}{\sum_{k\in S}d_k^{-p}},
\label{eq:idw}
\end{equation}
con $S$ los $k$ nodos más cercanos a $P$ y $d_i$ la distancia en la proyección
Origen-Nacional (EPSG:9377). Valores por defecto: $p=2$, $k=16$. Como los pesos se
normalizan, el factor de escala de la proyección se cancela, y la distancia proyectada es
equivalente a la distancia geodésica a la escala de unos pocos kilómetros. Si $P$ coincide
con un nodo (a menos de \SI{1}{mm}), ese nodo recibe peso~1. El IDW no reproduce gradientes:
el valor interpolado siempre queda dentro del rango de los nodos vecinos y tiende a
``aplanar'' el campo hacia su media local.

\subsection{Kriging ordinario}
Los pesos minimizan la varianza del error de predicción bajo la condición de
insesgamiento $\sum\lambda_i=1$:
\begin{equation}
\begin{pmatrix}\boldsymbol{\Gamma} & \mathbf{1}\\ \mathbf{1}^T & 0\end{pmatrix}
\begin{pmatrix}\boldsymbol{\lambda}\\ \mu\end{pmatrix}=
\begin{pmatrix}\boldsymbol{\gamma}_0\\ 1\end{pmatrix},
\qquad
\sigma^2_{KO}=\boldsymbol{\lambda}^T\boldsymbol{\gamma}_0+\mu,
\label{eq:ko}
\end{equation}
donde $\boldsymbol\Gamma$ contiene las semivarianzas entre los $n$ nodos vecinos y
$\boldsymbol\gamma_0$ las semivarianzas entre cada nodo y $P$ \citep{cressie1993}. El
semivariograma experimental se ajusta por mínimos cuadrados ponderados por el número de
pares, eligiendo entre los modelos esférico, exponencial, gaussiano y circular, igual que
en el Ejercicio~2. Se tomaron tres decisiones propias de este ejercicio:
\begin{enumerate}
    \item \textbf{Variograma local.} $\zeta$ tiene tendencias regionales de decenas de
    metros, así que un variograma global no representaría la variabilidad a la escala de la
    interpolación. Por eso el semivariograma se ajusta, para cada punto, con los nodos de una
    ventana de $2s\times2s$ nodos (semiancho por defecto $s=6$, es decir, 144 nodos,
    $\approx\pm\SI{55}{km}$).
    \item \textbf{$\gamma(0)=0$.} El efecto pepita es un salto en $h\to0^+$, no el valor en
    el origen. Así el kriging es exacto en los nodos y la diagonal de $\boldsymbol\Gamma$ es
    nula.
    \item \textbf{Sin muestreo aleatorio de pares}: se usan todos los pares de la ventana para
    que el ajuste sea reproducible.
\end{enumerate}
El sistema se resuelve con los $n=16$ vecinos más cercanos.

\subsection{Cobertura y no extrapolación}
Antes de aplicar cualquier método se comprueba que $P$ esté dentro del rectángulo de la
malla (límites incluidos). Si no lo está, no se calcula ningún valor y se informa
``fuera de cobertura''. Además, cada método valida su propio soporte: la bicúbica, por
ejemplo, no es aplicable en las celdas del borde exterior. Ningún resultado se sustituye en
silencio por el de otro método.

% =====================================================================
\section{Procedimiento}
\label{sec:procedimiento}
% =====================================================================

\subsection{Flujo de trabajo}
\begin{center}\small
\fbox{\parbox{0.92\linewidth}{\centering
\textbf{1.} Cargar QGeoidCOL2023.TXT $\;\rightarrow\;$
\textbf{2.} Validar y reconstruir la malla $\;\rightarrow\;$
\textbf{3.} Autoverificar contra \texttt{interp2} $\;\rightarrow\;$
\textbf{4.} Ingresar punto(s) $(\varphi,\lambda,h)$ $\;\rightarrow\;$
\textbf{5.} Control de cobertura $\;\rightarrow\;$
\textbf{6.} Interpolar con el método elegido (nodos, pesos, contribuciones) $\;\rightarrow\;$
\textbf{7.} $H^N=h-\zeta$ $\;\rightarrow\;$
\textbf{8.} Vista general y local $\;\rightarrow\;$
\textbf{9.} Comparación de los seis métodos $\;\rightarrow\;$
\textbf{10.} Validación cruzada y sensibilidad $\;\rightarrow\;$
\textbf{11.} Comparación con GeoCol2004 y EGM2008 $\;\rightarrow\;$
\textbf{12.} Exportación (CSV, TXT, PNG)
}}
\end{center}

\subsection{Arquitectura del código}
\begin{sloppypar}
La herramienta es un archivo único y autocontenido (\texttt{appQGeoidCOL2023.m},
Anexo~\ref{anx:codigo}) que no depende del Mapping Toolbox ni del Statistics Toolbox.
Reutiliza la arquitectura de la aplicación del Ejercicio~2: interfaz \texttt{uifigure}, estado en
\texttt{fig.UserData}, bitácora, proyección \texttt{proyectarOrigenNacional} y kriging
por mínimos cuadrados ponderados. Lo que cambia es el cargue de datos, la validación y los
métodos, que se adaptaron a una malla regular. El núcleo es
\texttt{interpolarPunto}: cada método (\texttt{pesosVecinoMasCercano},
\texttt{pesosTriangulos}, \texttt{pesosBilineal}, \texttt{pesosBicubica},
\texttt{pesosIDW}, \texttt{pesosKriging}) devuelve únicamente los \emph{índices de los
nodos y sus pesos}, y $\hat\zeta=\sum w_i\zeta_i$ se calcula en un solo lugar, igual para
todos los métodos. Así, la tabla de nodos, pesos y contribuciones y la verificación
$\sum w_i=1$ salen del mismo cálculo que produce el resultado y no de un cálculo aparte.
\end{sloppypar}

\begin{lstlisting}[language=Matlab, caption={Núcleo común: suma ponderada, suma de pesos y tabla de nodos (fragmento de \texttt{interpolarPunto}).}]
z = M.Z(sub2ind(size(M.Z), iLat, iLon));
if any(isnan(z(w ~= 0)))
    res.motivo = 'Soporte incompleto: falta en la malla al menos un nodo con peso no nulo.';
    return
end
res.zeta = sum(w .* zc);
res.sumaPesos = sum(w);
res.nodos = table((1:numel(w)).', M.lat(iLat), M.lon(iLon), z, w, w .* z, d, ...);
\end{lstlisting}

\begin{lstlisting}[language=Matlab, caption={Pesos de la interpolación lineal por triángulos con diagonal fija SW--NE (\texttt{pesosTriangulos}).}]
[i, j, u, v] = ubicarCelda(M, lat0, lon0);
if u >= v   % triangulo SW-SE-NE
    iLat = [i; i; i+1]; iLon = [j; j+1; j+1]; w = [1-u; u-v; v];
else        % triangulo SW-NE-NW
    iLat = [i; i+1; i+1]; iLon = [j; j+1; j]; w = [1-v; u; v-u];
end
\end{lstlisting}

\subsection{Parámetros y decisiones metodológicas}
\begin{table}[H]
\centering
\caption{Parámetros y decisiones de la implementación.}
\label{tab:parametros}
\small
\renewcommand{\arraystretch}{1.2}
\begin{tabular}{@{}p{3.6cm}p{11.6cm}@{}}
\toprule
Aspecto & Decisión y justificación \\
\midrule
Métrica & Los métodos sobre la malla operan con índices $(u,v)$. IDW y kriging usan
distancias en Origen-Nacional (EPSG:9377, GRS80, $k_0=0{,}9992$), en km. \\
Lineal por triángulos & Diagonal fija SW--NE (resultado reproducible). \\
Bicúbica & Keys, $a=-0{,}5$. No aplicable en la celda de borde (sin sustitución). \\
IDW & $p=2$, $k=16$ nodos más cercanos. \\
Kriging & Variograma local (semiancho de 6 nodos, 12 rezagos), $n=16$ vecinos,
$\gamma(0)=0$ y se reporta el \texttt{rcond} del sistema. \\
Cobertura & Rectángulo de la malla. Fuera de él no se calcula nada. \\
Precisión & Doble precisión en todo el cálculo. Los CSV se exportan sin redondear. \\
Redondeo de presentación & $\zeta$ y $H^N$ a \SI{0.001}{m} (\SI{1}{mm}, la resolución con que viene
$\zeta$ en el archivo); pesos a $10^{-6}$; contribuciones a $10^{-4}$\,m. \\
Reproducibilidad & Sin pasos aleatorios sin semilla. La validación cruzada y la autoverificación
usan \texttt{RandStream} con semilla fija, sin alterar el generador global de MATLAB. \\
\bottomrule
\end{tabular}
\end{table}

\subsection{Autoverificación de la implementación}
\label{sec:autoverif}
Al cargar la malla, la herramienta compara sus propias implementaciones con las
funciones de MATLAB en 300 puntos interiores aleatorios (semilla fija) y verifica la suma
de pesos de los seis métodos (Tabla~\ref{tab:autoverif}). Las diferencias son del orden
del redondeo de punto flotante, lo que confirma que las fórmulas de las
ecuaciones~\eqref{eq:bilineal} y~\eqref{eq:keys} están bien implementadas, incluida la
orientación de la malla.

\begin{table}[H]
\centering
\caption{Autoverificación: diferencias máximas frente a \texttt{interp2} y error máximo de la suma de pesos.}
\label{tab:autoverif}
\begin{tabular}{@{}lll@{}}
\toprule
Método & $\max|\hat\zeta-\hat\zeta_{\texttt{interp2}}|$ (m) & $\max|1-\sum w_i|$ \\
\midrule
Vecino más cercano & 0 (\texttt{'nearest'}) & 0 \\
Lineal por triángulos & --- & \num{1.1e-16} \\
Bilineal & \num{7.1e-15} (\texttt{'linear'}) & \num{2.2e-16} \\
Bicúbica & \num{6.8e-14} (\texttt{'cubic'}) & \num{2.3e-15} \\
IDW & --- & \num{5.6e-16} \\
Kriging ordinario & --- & \num{4.4e-16} \\
\bottomrule
\end{tabular}
\end{table}

\subsection{Validación cruzada por decimación}
\label{sec:cvmetodo}
En el Ejercicio~2 se mostró que el \emph{leave-one-out} clásico no sirve para métodos
exactos de malla: al retirar un nodo y consultar justo en su posición, el peso cae por
completo sobre el nodo ausente. Para este ejercicio se diseñó una validación adaptada a
mallas: se construye una malla \emph{decimada} con los nodos de índice impar (paso de
$10'$) y con ella se predicen 400 nodos retirados, elegidos al azar con semilla 2023. Esos
nodos quedan en centros de celda y en puntos medios de aristas de la malla decimada. El
error es $e=\hat\zeta-\zeta_{\text{archivo}}$. Como el paso de prueba es el doble del real,
los errores son una \textbf{cota conservadora} del error al paso de $5'$ y sirven para
comparar los métodos entre sí, no como exactitud absoluta del modelo.

% =====================================================================
\section{Interfaz gráfica}
\label{sec:interfaz}
% =====================================================================
La aplicación se ejecuta con \texttt{appQGeoidCOL2023} y carga automáticamente
\texttt{QGeoidCOL2023.txt} si está en la misma carpeta; si no, se usa el botón ``Cargar''.
El panel izquierdo contiene la información del modelo, el selector de método, los
parámetros de IDW y kriging, el cuadro de puntos (formato \texttt{[Nombre:] lat, lon [, h]})
y los botones de cálculo, comparación, validación y exportación, además de la bitácora. El panel
derecho tiene cinco pestañas: malla, consulta y vista local, comparación de métodos,
validación cruzada y comparación entre modelos (Figuras~\ref{fig:app1}--\ref{fig:app3}).
En la pestaña de modelos, la tabla de entrada admite la ondulación de GeoCol2004, la de EGM2008 y
la altura nivelada $H$ (digitadas o importadas de un CSV, como \texttt{valores\_control\_P1\_P5.csv}) y calcula
las diferencias, las alturas de cada modelo y su diferencia con la nivelación. También existe un modo por lotes
sin interfaz, \texttt{appQGeoidCOL2023('lote')}, y uno de sensibilidad,
\texttt{appQGeoidCOL2023('sensibilidad')}, que reproducen los resultados numéricos de este informe;
\texttt{generarResultadosInforme.m} regenera además las figuras y los archivos de datos (Anexo~\ref{anx:archivos}).

\begin{figure}[H]
\centering
\includegraphics[width=\linewidth]{figuras/app_01_carga.png}
\caption{Interfaz recién cargada la malla: vista general de $\zeta$, histograma y bitácora con
el reporte de validación.}
\label{fig:app1}
\end{figure}

\begin{figure}[H]
\centering
\includegraphics[width=\linewidth]{figuras/app_02_consulta_kriging.png}
\caption{Pestaña de consulta con kriging ordinario en P4: vista local con nodos, pesos
(en azul los negativos) y geometría de soporte; resumen con $\zeta$, $h$, $H^N$, suma de
pesos y parámetros del variograma; tabla de nodos con su contribución.}
\label{fig:app2}
\end{figure}

\begin{figure}[H]
\centering
\includegraphics[width=\linewidth]{figuras/app_04_comparacion.png}
\caption{Pestaña de comparación de los seis métodos en los cinco puntos, con la bilineal
como referencia.}
\label{fig:app3}
\end{figure}

\subsection{Diagrama de flujo}
\begin{figure}[H]
\centering
\adjustbox{max width=0.98\linewidth}{%
\begin{tikzpicture}[
    font=\scriptsize,
    >={Stealth[length=2mm]},
    line/.style={->,line width=0.55pt},
    start/.style={ellipse,draw,fill=gray!15,minimum width=28mm,minimum height=7mm,align=center},
    io/.style={trapezium,trapezium left angle=70,trapezium right angle=110,draw,fill=orange!12,
        minimum width=30mm,minimum height=8mm,align=center},
    proc/.style={rectangle,rounded corners=2pt,draw,fill=gray!8,minimum width=34mm,minimum height=8mm,align=center},
    decision/.style={diamond,draw,fill=yellow!25,aspect=2.2,inner sep=1pt,align=center},
    method/.style={rectangle,rounded corners=2pt,draw,fill=blue!10,minimum width=22mm,minimum height=9mm,align=center},
    outp/.style={rectangle,rounded corners=2pt,draw,fill=green!12,minimum width=34mm,minimum height=8mm,align=center},
    err/.style={rectangle,draw,fill=red!12,minimum width=24mm,minimum height=7mm,align=center},
]
\node[start] (inicio) {Inicio\\\texttt{appQGeoidCOL2023}};
\node[io,below=6mm of inicio] (cargar) {Cargar\\\texttt{QGeoidCOL2023.TXT}};
\node[proc,below=6mm of cargar] (validar) {\texttt{cargarYValidarMalla}\\numéricos, finitos, rango,\\duplicados, regularidad, faltantes};
\node[decision,below=7mm of validar] (dec1) {¿Malla\\regular?};
\node[err,right=14mm of dec1] (err1) {Error: malla\\no regular};
\node[proc,below=7mm of dec1] (auto) {Autoverificación\\vs \texttt{interp2}, $\sum w=1$};
\node[io,below=6mm of auto] (punto) {Punto(s) $(\varphi,\lambda,h)$};
\node[decision,below=7mm of punto] (dec2) {¿Dentro de\\la malla?};
\node[err,right=14mm of dec2] (err2) {Fuera de\\cobertura:\\sin valor};
\node[decision,below=8mm of dec2] (met) {Método};

\node[method,below=12mm of met,xshift=-62mm] (m1) {Vecino\\más cercano};
\node[method,right=3mm of m1] (m2) {Lineal por\\triángulos};
\node[method,right=3mm of m2] (m3) {Bilineal};
\node[method,right=3mm of m3] (m4) {Bicúbica\\(Keys)};
\node[method,right=3mm of m4] (m5) {IDW\\($p$, $k$)};
\node[method,right=3mm of m5] (m6) {Kriging\\ordinario};
\node[err,below=5mm of m4] (err3) {¿Soporte 4$\times$4?\\No: no aplicable};

\node[proc,below=20mm of m3,xshift=12mm] (nucleo) {$\hat\zeta=\sum w_i\zeta_i$, verificación $\sum w_i=1$\\tabla de nodos, pesos y contribuciones};
\node[proc,below=6mm of nucleo] (hn) {$H^N=h-\hat\zeta$ (si hay $h$)};
\node[outp,below=6mm of hn] (out) {Vista general y local, tablas,\\comparación, validación, exportación};

\draw[line] (inicio) -- (cargar);
\draw[line] (cargar) -- (validar);
\draw[line] (validar) -- (dec1);
\draw[line] (dec1.east) -- node[above]{No} (err1.west);
\draw[line] (dec1.south) -- node[left]{Sí} (auto.north);
\draw[line] (auto) -- (punto);
\draw[line] (punto) -- (dec2);
\draw[line] (dec2.east) -- node[above]{No} (err2.west);
\draw[line] (dec2.south) -- node[left]{Sí} (met.north);
\coordinate (rep) at ($(met.south)+(0,-5mm)$);
\draw (met.south) -- (rep);
\foreach \m in {m1,m2,m3,m4,m5,m6} {\draw[line] (rep) -| (\m.north);}
\draw[line,dashed] (m4.south) -- (err3.north);
\coordinate (conv) at ($(nucleo.north)+(0,5mm)$);
\foreach \m in {m1,m2,m3,m5,m6} {\draw (\m.south) |- (conv);}
\draw (err3.south) |- (conv);
\draw[line] (conv) -- (nucleo.north);
\draw[line] (nucleo) -- (hn);
\draw[line] (hn) -- (out);
\end{tikzpicture}}
\caption{Diagrama de flujo de \texttt{appQGeoidCOL2023}.}
\label{fig:diagrama}
\end{figure}

% =====================================================================
\section{Resultados}
\label{sec:resultados}
% =====================================================================

\subsection{Anomalía de altura y altura normal en los puntos de control}
La Tabla~\ref{tab:zeta} reúne $\hat\zeta$ por los seis métodos en los cinco puntos (los
cinco están dentro de la cobertura y ninguno cae en una celda de borde, así que los seis
métodos son aplicables en todos). La Tabla~\ref{tab:hn} presenta $h$, $\zeta$ y $H^N$
para el método de referencia (bilineal).

\begin{table}[H]
\centering
\caption{Anomalía de altura $\hat\zeta$ (m) por método. Calculada en doble precisión y
presentada con cuatro decimales para hacer visibles las diferencias submilimétricas.}
\label{tab:zeta}
\small
\begin{tabular}{@{}lrrrrrr@{}}
\toprule
Punto & Vec. cercano & Triángulos & Bilineal & Bicúbica & IDW & Kriging \\
\midrule
P1 & \num{-13.4480} & \num{-13.2702} & \num{-13.2822} & \num{-13.2931} & \num{-13.3012} & \num{-13.2927} \\
P2 & \num{2.3150} & \num{2.5327} & \num{2.5333} & \num{2.5237} & \num{2.5075} & \num{2.5234} \\
P3 & \num{11.8490} & \num{11.7409} & \num{11.7637} & \num{11.8233} & \num{11.7469} & \num{11.7404} \\
P4 & \num{23.3320} & \num{23.3668} & \num{23.3603} & \num{23.3722} & \num{23.3247} & \num{23.3564} \\
P5 & \num{0.5750} & \num{0.8196} & \num{0.8260} & \num{0.8057} & \num{0.8076} & \num{0.8048} \\
\bottomrule
\end{tabular}
\end{table}

\begin{table}[H]
\centering
\caption{Altura normal $H^N=h-\zeta$ con la interpolación bilineal (redondeo a \SI{1}{mm}).}
\label{tab:hn}
\begin{tabular}{@{}lrrr@{}}
\toprule
Punto & $h$ (m) & $\zeta$ (m) & $H^N$ (m) \\
\midrule
P1 & \num{35.134} & \num{-13.282} & \num{48.416} \\
P2 & \num{98.266} & \num{2.533} & \num{95.733} \\
P3 & \num{3228.966} & \num{11.764} & \num{3217.202} \\
P4 & \num{2725.381} & \num{23.360} & \num{2702.021} \\
P5 & \num{185.792} & \num{0.826} & \num{184.966} \\
\bottomrule
\end{tabular}
\end{table}

\subsection{Nodos, pesos y contribuciones (P4)}
Se detalla el punto P4 (Bogotá, $u=0{,}8455$, $v=0{,}7796$), que cae en una zona de
gradiente fuerte de $\zeta$. Como $u\ge v$, el triángulo que lo contiene es SW--SE--NE.
Las Tablas~\ref{tab:nodos_bil} y~\ref{tab:nodos_tri} muestran los nodos de la bilineal y
de la lineal por triángulos, y la Figura~\ref{fig:locales} muestra la vista local de los
seis métodos. Las tablas completas de IDW y kriging (16 nodos cada una) están en el
Anexo~\ref{anx:tablas}.

\begin{table}[H]
\centering
\caption{P4, bilineal: nodos, pesos y contribuciones. $\sum w_i=1{,}000000000000000$; $\hat\zeta=\SI{23.3603}{m}$.}
\label{tab:nodos_bil}
\small
\begin{tabular}{@{}lrrrrrr@{}}
\toprule
Nodo & $\varphi$ (grados) & $\lambda$ (grados) & $\zeta_i$ (m) & $w_i$ & $w_i\zeta_i$ (m) & $d$ (km) \\
\midrule
SW & \num{4.550000} & \num{-74.133333} & \num{23.388} & \num{0.034043} & \num{0.7962} & \num{10.612} \\
SE & \num{4.550000} & \num{-74.050000} & \num{23.729} & \num{0.186318} & \num{4.4211} & \num{7.321} \\
NE & \num{4.633333} & \num{-74.050000} & \num{23.332} & \num{0.659196} & \num{15.3804} & \num{2.481} \\
NW & \num{4.633333} & \num{-74.133333} & \num{22.937} & \num{0.120443} & \num{2.7626} & \num{8.073} \\
\midrule
$\sum$ & & & & \num{1.000000} & \num{23.3603} & \\
\bottomrule
\end{tabular}
\end{table}

\begin{table}[H]
\centering
\caption{P4, lineal por triángulos (SW--SE--NE): pesos $(1-u,\,u-v,\,v)$. $\hat\zeta=\SI{23.3668}{m}$.}
\label{tab:nodos_tri}
\small
\begin{tabular}{@{}lrrrrrr@{}}
\toprule
Nodo & $\varphi$ (grados) & $\lambda$ (grados) & $\zeta_i$ (m) & $w_i$ & $w_i\zeta_i$ (m) & $d$ (km) \\
\midrule
SW & \num{4.550000} & \num{-74.133333} & \num{23.388} & \num{0.154486} & \num{3.6131} & \num{10.612} \\
SE & \num{4.550000} & \num{-74.050000} & \num{23.729} & \num{0.065875} & \num{1.5631} & \num{7.321} \\
NE & \num{4.633333} & \num{-74.050000} & \num{23.332} & \num{0.779639} & \num{18.1905} & \num{2.481} \\
\midrule
$\sum$ & & & & \num{1.000000} & \num{23.3668} & \\
\bottomrule
\end{tabular}
\end{table}

En P4, el vecino más cercano toma el nodo NE ($w=1$, a \SI{2.481}{km}). La bicúbica usa
16 nodos, con el mayor peso en el nodo NE ($0{,}846$) y ocho pesos negativos en la corona
exterior. El IDW concentra el 63\,\% del peso en el nodo más cercano, y el kriging lo
reparte más ($0{,}248$ en el más cercano) y asigna pesos negativos a los seis nodos más
lejanos. El variograma local que eligió el kriging en P4 fue gaussiano, con pepita de
\SI{0.064}{m^2}, meseta parcial de \SI{4.15}{m^2}, rango de \SI{40.8}{km} y
$\sigma_{KO}=\SI{0.283}{m}$.

\begin{figure}[H]
\centering
\begin{subfigure}[t]{0.32\linewidth}\centering
\includegraphics[width=\linewidth]{figuras/local_P4_vmc.png}\caption{Vecino más cercano}\end{subfigure}
\hfill
\begin{subfigure}[t]{0.32\linewidth}\centering
\includegraphics[width=\linewidth]{figuras/local_P4_tri.png}\caption{Lineal por triángulos}\end{subfigure}
\hfill
\begin{subfigure}[t]{0.32\linewidth}\centering
\includegraphics[width=\linewidth]{figuras/local_P4_bil.png}\caption{Bilineal}\end{subfigure}

\vspace{0.3cm}
\begin{subfigure}[t]{0.32\linewidth}\centering
\includegraphics[width=\linewidth]{figuras/local_P4_bic.png}\caption{Bicúbica}\end{subfigure}
\hfill
\begin{subfigure}[t]{0.32\linewidth}\centering
\includegraphics[width=\linewidth]{figuras/local_P4_idw.png}\caption{IDW ($p=2$, $k=16$)}\end{subfigure}
\hfill
\begin{subfigure}[t]{0.32\linewidth}\centering
\includegraphics[width=\linewidth]{figuras/local_P4_ko.png}\caption{Kriging ordinario}\end{subfigure}
\caption{Vista local en P4: campo interpolado con cada método ($\pm2{,}5$ nodos), nodos
de la malla (puntos negros), nodos usados (tamaño proporcional a $|w_i|$; rojo $w_i>0$,
azul $w_i<0$, etiqueta con el peso) y geometría de soporte (blanco: triángulo, celda,
ventana $4\times4$ o líneas radiales; discontinua: celda que contiene al punto).}
\label{fig:locales}
\end{figure}

\subsection{Diferencias entre métodos (m, cm, mm)}
Se tomó la interpolación bilineal como referencia, por ser el método habitual de consulta
de mallas geoidales, y no porque se considere un valor verdadero. La Tabla~\ref{tab:difmm}
presenta $\hat\zeta_{\text{método}}-\hat\zeta_{\text{bilineal}}$ y la Tabla~\ref{tab:dispersion}
resume la dispersión en las tres unidades.

\begin{table}[H]
\centering
\caption{Diferencia respecto a la bilineal, $\hat\zeta_{\text{método}}-\hat\zeta_{\text{bilineal}}$, en mm.}
\label{tab:difmm}
\begin{tabular}{@{}lrrrrr@{}}
\toprule
Punto & Vec. cercano & Triángulos & Bicúbica & IDW & Kriging \\
\midrule
P1 & \num{-165.8} & \num{12.0} & \num{-10.9} & \num{-19.0} & \num{-10.5} \\
P2 & \num{-218.3} & \num{-0.6} & \num{-9.6} & \num{-25.8} & \num{-9.9} \\
P3 & \num{85.3} & \num{-22.8} & \num{59.6} & \num{-16.8} & \num{-23.3} \\
P4 & \num{-28.3} & \num{6.5} & \num{11.9} & \num{-35.6} & \num{-3.9} \\
P5 & \num{-251.0} & \num{-6.3} & \num{-20.3} & \num{-18.3} & \num{-21.2} \\
\bottomrule
\end{tabular}
\end{table}

\begin{figure}[H]
\centering
\includegraphics[width=\linewidth]{figuras/1.png}
\caption{Diferencia respecto a la interpolación bilineal en la anomalía de altura $\hat\zeta$ (mm) por método}
\label{fig:app1}
\end{figure}

\begin{table}[H]
\centering
\caption{Dispersión (máximo menos mínimo) de $\hat\zeta$ entre métodos en cada punto.}
\label{tab:dispersion}
\begin{tabular}{@{}lrrr@{\hspace{1.2cm}}rrr@{}}
\toprule
& \multicolumn{3}{c}{Seis métodos} & \multicolumn{3}{c}{Sin el vecino más cercano} \\
\cmidrule(lr){2-4}\cmidrule(l){5-7}
Punto & m & cm & mm & m & cm & mm \\
\midrule
P1 & \num{0.178} & \num{17.8} & \num{177.8} & \num{0.031} & \num{3.1} & \num{31.0} \\
P2 & \num{0.218} & \num{21.8} & \num{218.3} & \num{0.026} & \num{2.6} & \num{25.8} \\
P3 & \num{0.109} & \num{10.9} & \num{108.6} & \num{0.083} & \num{8.3} & \num{82.9} \\
P4 & \num{0.048} & \num{4.8} & \num{47.6} & \num{0.048} & \num{4.8} & \num{47.6} \\
P5 & \num{0.251} & \num{25.1} & \num{251.0} & \num{0.021} & \num{2.1} & \num{21.2} \\
\bottomrule
\end{tabular}
\end{table}

\subsection{Validación cruzada}
\begin{table}[H]
\centering
\caption{Validación cruzada por decimación (malla de prueba con paso de $10'$, 400 nodos
de control, semilla 2023). Errores en mm.}
\label{tab:cv}
\begin{tabular}{@{}lrrrrrr@{}}
\toprule
Método & Evaluados & Sesgo & RMSE & MAE & P95 $|e|$ & Máx.\ $|e|$ \\
\midrule
Vecino más cercano & 400 & \num{88.5} & \num{480.7} & \num{326.3} & \num{997.0} & \num{2532.0} \\
Lineal por triángulos & 400 & \num{-1.5} & \num{88.7} & \num{49.8} & \num{184.0} & \num{611.5} \\
Bilineal & 400 & \num{-0.4} & \num{86.6} & \num{48.3} & \num{187.2} & \num{543.7} \\
Bicúbica & 379$^{*}$ & \num{-0.9} & \num{58.0} & \num{29.5} & \num{141.4} & \num{313.2} \\
IDW ($p=2$, $k=16$) & 400 & \num{6.6} & \num{212.3} & \num{133.4} & \num{449.3} & \num{959.4} \\
Kriging ordinario & 400 & \num{-6.3} & \num{135.9} & \num{64.3} & \num{315.6} & \num{856.4} \\
\bottomrule
\end{tabular}

\vspace{0.1cm}
\footnotesize $^{*}$21 controles quedaron en celdas de borde de la malla decimada, donde la
bicúbica no es aplicable (no se sustituyeron).
\end{table}

Para medir qué tan bien interpola cada método se tiene que el valor real se conoce,
la diferencia es directamente el error. El \textbf{sesgo}
es el error promedio con signo (cerca de cero = sin desviación sistemática); el
\textbf{RMSE} es el error típico y penaliza las fallas grandes; el \textbf{MAE}
es el error típico sin elevar al cuadrado; el \textbf{P95} es el valor que no
supera el $95\%$ de los errores; y el \textbf{máximo} es el peor caso, de los 400 nodos de prueba, el método que falló por mayor diferencia fue vecino más cercano con $2532 mm$, o sea unos $2.5m$

\vspace{0.4cm}
El \textbf{diagrama de cajas} (izquierda) muestra cómo se reparten los $400$
errores de cada método: la caja encierra la mitad central, la línea interna es
el error típico y los puntos sueltos son casos atípicos. Las \textbf{barras}
(derecha) comparan los cuatro indicadores por método en escala logarítmica,
necesaria porque los errores van de decenas a miles de milímetros.

\vspace{0.4cm}
La interpolación por \textbf{vecino más cercano} presenta el peor (RMSE
$\approx 481$~mm y sesgo notorio), mientras que la \textbf{bicúbica} es la que presenta mejor
(RMSE $\approx 58$~mm), seguida de cerca por la \textbf{lineal por triángulos} y
la \textbf{bilineal}; el \textbf{IDW} y el \textbf{kriging ordinario} quedan en
un punto intermedio.

\begin{figure}[H]
\centering
\begin{subfigure}[t]{0.44\linewidth}\centering
\includegraphics[width=\linewidth]{figuras/cv_boxplot.png}\caption{Distribución de errores.}\end{subfigure}
\hfill
\begin{subfigure}[t]{0.44\linewidth}\centering
\includegraphics[width=\linewidth]{figuras/cv_indicadores.png}\caption{Indicadores (escala logarítmica).}\end{subfigure}
\caption{Validación cruzada por decimación.}
\label{fig:cv}
\end{figure}

\subsection{Sensibilidad de parámetros}
\label{sec:sensibilidad}
Con los mismos 400 nodos de control se evaluaron combinaciones de los parámetros de IDW
($p$, $k$) y de kriging ($n$, semiancho de la ventana) (Tabla~\ref{tab:sensibilidad}).

\begin{table}[H]
\centering
\caption{Sensibilidad de parámetros: RMSE (mm) de la validación cruzada por decimación
(mismos 400 controles). En negrita, la configuración por defecto de la herramienta.}
\label{tab:sensibilidad}
\begin{tabular}{@{}lrrrr@{\hspace{1.2cm}}lrrr@{}}
\toprule
\multicolumn{5}{c}{IDW} & \multicolumn{4}{c}{Kriging ordinario} \\
\cmidrule(r){1-5}\cmidrule(l){6-9}
$p$ \textbackslash{} $k$ & 4 & 8 & 16 & 32 & $n$ \textbackslash{} semiancho & 4 & 6 & 9 \\
\midrule
1 & \num{240.8} & \num{194.4} & \num{323.0} & \num{530.2} & 8  & \num{94.6} & \num{139.0} & \num{172.3} \\
2 & \num{150.4} & \num{143.7} & \textbf{\num{212.3}} & \num{310.8} & 16 & \num{85.7} & \textbf{\num{135.9}} & \num{222.0} \\
3 & \num{106.9} & \num{111.7} & \num{139.2} & \num{169.2} & 32 & \num{103.9} & \num{147.0} & \num{226.2} \\
4 & \num{92.2} & \num{96.8} & \num{106.4} & \num{113.4} & & & & \\
\bottomrule
\end{tabular}

\vspace{0.1cm}
\footnotesize Referencias en la misma prueba: bilineal \SI{86.6}{mm}; bicúbica \SI{58.0}{mm}.
\end{table}

\subsection{Modelo paramétrico de la discrepancia con la Red Nacional de Nivelación (Python)}
\label{sec:parametrico}

Como complemento a la interpolación de la malla se desarrolló en Python una segunda
herramienta. Con ella se evalúa qué tan bien describen los modelos de geoide las alturas de
la Red Nacional de Nivelación, usando los puntos del archivo \texttt{coords\_QGEOID2023.txt}
(IGAC). Las mallas analizadas son QGeoidCOL2023 y EGM2008. Esta última se generó en el
servicio ICGEM \citep{ince2019} para la región de
Colombia\footnote{\url{https://icgem.gfz.de/calc_res/d533101e8204f4397361b8d33267c6ed173ceea602f93b06f39678a51579bc4c}}.
El procedimiento sigue el planteamiento de \citet{guimaraes2025} para conectar los sistemas
verticales locales de Sudamérica con el IHRF y tiene cuatro etapas: homogeneizar el sistema de
mareas, depurar los datos, interpolar el geoide con los seis métodos de la
Sección~\ref{sec:marco} y ajustar un modelo paramétrico a la discrepancia.

\subsubsection{Discrepancia entre alturas y geoide}
En cada punto nivelado se conocen la altura elipsoidal $h$ y la altura sobre el nivel del mar
$H$. La ondulación $N$ se interpola de la malla en la posición $(\varphi,\lambda)$ del punto.
Según la ecuación~\eqref{eq:alturas}, si los tres valores fueran exactos se cumpliría $h=H+N$.
En la práctica queda una diferencia, la \textbf{discrepancia}:
\begin{equation}
l = h - H - N .
\label{eq:discrepancia}
\end{equation}
Si $l$ fuera igual en todos los puntos, el nivel de referencia de la red (el datum de
Buenaventura) estaría desplazado una cantidad constante respecto al geoide del modelo. Si
$l$ cambia de un punto a otro, hay inclinaciones de la red, errores de nivelación o
imprecisiones del modelo. En la red colombiana $H$ es una altura normal-ortométrica, así que,
tanto con QGeoidCOL2023 (que entrega $\zeta$) como con EGM2008 (que entrega $N$), la
discrepancia incluye además la diferencia conceptual de la ecuación~\eqref{eq:nzeta}.

\subsubsection{Homogeneización al sistema de mareas \emph{zero tide}}
La atracción permanente del Sol y la Luna deforma la Tierra y el geoide. Según cómo se trate
ese efecto existen tres convenciones: \emph{tide free}, \emph{mean tide} y \emph{zero tide}. Para
que la ecuación~\eqref{eq:discrepancia} tenga sentido, $h$, $H$ y $N$ deben estar en la misma
convención, y aquí se adopta \emph{zero tide}, que es la del IHRS. El usuario declara el sistema
de cada dato y la herramienta solo corrige lo que no está en \emph{zero tide}. En los datos de
este ejercicio las alturas elipsoidales vienen en \emph{tide free}, las niveladas en \emph{mean
tide}, QGeoidCOL2023 en \emph{zero tide} y EGM2008 en \emph{tide free}, como lo indica el
encabezado del archivo de ICGEM. Para las alturas se usan las expresiones de \citet{guimaraes2025}:
\begin{align}
H_{ZT} &= H_{MT} - 0{,}198\left(\tfrac{3}{2}\sin^2\varphi-\tfrac{1}{2}\right), \label{eq:mareaH}\\
h_{ZT} &= h_{TF} + 0{,}0602 - 0{,}1790\,\sin^2\varphi - 0{,}0018\,\sin^4\varphi, \label{eq:mareah}
\end{align}
en metros. El artículo no da una expresión para las mallas; para ellas se usó la que se
deduce de la ecuación~\eqref{eq:mareaH} con el número de Love $k=0{,}3$:
\begin{equation}
N_{ZT} = N_{TF} + k\left[-0{,}198\left(\tfrac{3}{2}\sin^2\varphi-\tfrac{1}{2}\right)\right],
\qquad k=0{,}3 .
\label{eq:mareaN}
\end{equation}
En el área de los puntos ($\varphi$ entre $1{,}3^\circ$ y $11{,}5^\circ$), $h$ aumenta entre
\SI{5.3}{cm} y \SI{6.0}{cm} y $H$ entre \SI{8.7}{cm} y \SI{9.9}{cm}. Por eso la discrepancia
disminuye entre \SI{3.4}{cm} y \SI{3.9}{cm}. La malla EGM2008 sube entre \SI{2.6}{cm} y \SI{3.0}{cm} y
se exporta convertida, con el mismo formato del archivo original.

\subsubsection{Depuración de los datos}
Antes del ajuste se revisan los puntos y se apartan los que tienen errores evidentes
(Tabla~\ref{tab:depuracion}). Ningún punto se descarta sin registro: todos quedan en
\texttt{puntos\_rechazados.xlsx} con la regla que incumplieron. La regla de error grueso
calcula la discrepancia de cada punto con interpolación bilineal y la compara con la del
resto mediante la mediana y la desviación absoluta mediana (MAD). Ninguno de los dos
estadísticos se deja arrastrar por los propios valores atípicos:
\begin{equation}
\left|\,l - \operatorname{med}(l)\,\right| > c\cdot 1{,}4826\cdot\operatorname{MAD}(l),
\qquad c=5 \text{ por defecto}.
\label{eq:grueso}
\end{equation}
El factor $1{,}4826$ hace que $1{,}4826\cdot\mathrm{MAD}$ sea comparable con una desviación
estándar, y un umbral $c$ mayor aparta menos puntos.

\begin{table}[H]
\centering
\caption{Reglas de depuración de los puntos de la Red Nacional de Nivelación.}
\label{tab:depuracion}
\small
\renewcommand{\arraystretch}{1.15}
\begin{tabular}{@{}lp{12.4cm}@{}}
\toprule
Regla & Se aparta el punto cuando \\
\midrule
R1 & falta un dato o un valor no es numérico. \\
R2 & la latitud o la longitud están fuera de rango. \\
R3 & el identificador o las coordenadas están repetidos. \\
R4 & $h=0$ con $H>\SI{1}{m}$ (altura elipsoidal faltante). \\
R5 & $|h-H|>\SI{150}{m}$ (ninguna ondulación alcanza ese valor). \\
R6 & su discrepancia cumple la condición~\eqref{eq:grueso} en alguna malla (error grueso). \\
R7 & queda fuera de la cobertura de alguna malla. \\
\bottomrule
\end{tabular}
\end{table}

\subsubsection{Interpolación y semivariograma del kriging}
El valor de $N$ en cada punto se obtiene con los seis métodos de la Sección~\ref{sec:marco}, con
las mismas implementaciones de la aplicación de MATLAB. En el kriging ordinario, el
semivariograma experimental de una ventana local se ajusta a los modelos esférico,
exponencial, gaussiano y circular, y se elige el de menor suma de cuadrados ponderada de los
residuos:
\begin{equation}
\mathrm{SSR_w}=\sum_{k} n_k\left[\hat\gamma(h_k)-\gamma(h_k;\theta)\right]^2,
\label{eq:ssrw}
\end{equation}
donde $n_k$ es el número de pares de nodos del rezago $k$. Así, un rezago calculado con más
pares, y por tanto más confiable, pesa más en el ajuste (Figura~\ref{fig:semivario}). Los
puntos que caen en la misma celda de la malla comparten ventana, de modo que se ajusta un
semivariograma por celda y no por punto: 
644 con QGeoidCOL2023 y 
1528 con EGM2008.
Para el IDW, el exponente $p$ puede sugerirse con la validación cruzada por decimación de la
Sección~\ref{sec:cvmetodo}, eligiendo el valor de menor RMSE entre $p=0{,}5$ y $p=6$. En
las dos mallas el valor sugerido fue $p=6$, y es el que se usó.

\begin{figure}[H]
\centering
\includegraphics[width=0.99\linewidth]{figuras/semivariokriging.png}
\caption{Semivariograma local del kriging ordinario en la herramienta de Python: valores
experimentales (tamaño proporcional al número de pares), los cuatro modelos ajustados y el
elegido por menor $\mathrm{SSR_w}$.}
\label{fig:semivario}
\end{figure}

\subsubsection{Modelo paramétrico y selección del método}
La parte sistemática de la discrepancia se representa con el modelo de \citet{guimaraes2025},
que supone un desplazamiento vertical común más una inclinación norte--sur y otra este--oeste:
\begin{equation}
l_P = \delta H + x_1\,(\varphi_P-\varphi_0) + x_2\,(\lambda_P-\lambda_0)\cos\varphi_P + v_P .
\label{eq:modelo18}
\end{equation}
En el modelo, $\delta H$ es el \textbf{desplazamiento vertical medio} del datum respecto al
geoide del modelo (m) y $x_1$ y $x_2$ son las inclinaciones norte--sur y este--oeste (m por
grado). $(\varphi_0,\lambda_0)$ es la posición media de los puntos y $v_P$ el residuo. El
factor $\cos\varphi_P$ hace comparables los grados de longitud con los de latitud. En forma
matricial, $\mathbf{l}=\mathbf{A}\mathbf{x}+\mathbf{v}$ con $\mathbf{x}=(\delta H,\,x_1,\,x_2)^T$, y el
sistema se resuelve por mínimos cuadrados con todos los puntos con el mismo peso
($\mathbf{P}=\mathbf{I}$):
\begin{equation}
\hat{\mathbf{x}}=(\mathbf{A}^T\mathbf{P}\mathbf{A})^{-1}\mathbf{A}^T\mathbf{P}\,\mathbf{l},
\qquad
\hat\sigma_0^2=\frac{\mathbf{v}^T\mathbf{v}}{n-u},
\qquad
\sigma_{\hat x_k}=\hat\sigma_0\sqrt{\left[(\mathbf{A}^T\mathbf{A})^{-1}\right]_{kk}} ,
\label{eq:mmcc}
\end{equation}
con $n$ el número de puntos y $u$ el de parámetros. El modelo se ajusta seis veces, una con
el $N$ de cada método y siempre con los mismos puntos. Para comparar los ajustes se usa el
coeficiente de determinación ajustado, que indica qué fracción de la variación de $l$
explica el modelo y corrige el $R^2$ según el número de parámetros:
\begin{equation}
R^2 = 1-\frac{\sum v_P^2}{\sum (l_P-\bar l)^2},
\qquad
R^2_{aj} = 1-(1-R^2)\,\frac{n-1}{n-u}.
\label{eq:r2aj}
\end{equation}
El mejor método es el de mayor $R^2_{aj}$; como criterio adicional se reporta $\hat\sigma_0$.
La herramienta también permite usar la superficie de corrección de cuatro parámetros,
$l_P=x_0+x_1\cos\varphi_P\cos\lambda_P+x_2\cos\varphi_P\sin\lambda_P+x_3\sin\varphi_P$. Es más
flexible, pero sus coeficientes no tienen una interpretación física individual.

\subsubsection{Estructura del código}
La implementación tiene dos archivos que deben estar en la misma carpeta y que solo usan la
biblioteca estándar de Python ($\geq 3.8$), sin \texttt{numpy} ni otras dependencias:
\begin{itemize}
    \item \texttt{parametrico1.py}: \textbf{núcleo de cálculo}. Hace todos los cálculos y
    también se puede ejecutar solo, desde la consola, en modo interactivo o con opciones de
    línea de comandos.
    \item \texttt{parametrico2.py}: \textbf{interfaz gráfica} en \texttt{tkinter}. Importa el
    núcleo (\texttt{import parametrico1 as mp}) y no repite ningún cálculo, por lo que las dos
    formas de uso dan resultados idénticos.
\end{itemize}
El núcleo está organizado por etapas (Tabla~\ref{tab:estructura_py}). Igual que en la
aplicación de MATLAB, cada método devuelve solo los índices de los nodos y sus pesos.
$\hat N=\sum w_iN_i$ se calcula en una función común (\texttt{interpolar}), que además
verifica la suma de pesos y rechaza el resultado si falta un nodo con peso no nulo.

\begin{table}[H]
\centering
\caption{Organización de \texttt{parametrico1.py} por etapas.}
\label{tab:estructura_py}
\small
\renewcommand{\arraystretch}{1.15}
\begin{tabular}{@{}p{3.4cm}p{11.8cm}@{}}
\toprule
Etapa & Funciones principales \\
\midrule
Sistema de mareas & \texttt{correccion\_h}, \texttt{correccion\_H}, \texttt{correccion\_N}
(ecuaciones~\eqref{eq:mareaH}--\eqref{eq:mareaN}); \texttt{aplicar\_marea\_malla},
\texttt{exportar\_malla\_zero\_tide}. \\
Lectura y depuración & \texttt{leer\_tabla}, \texttt{leer\_puntos} (R1--R5),
\texttt{cargar\_malla} (malla \emph{lat lon N} o formato ICGEM), \texttt{depurar\_contra\_mallas} (R6--R7). \\
Interpolación & \texttt{interpolar} (núcleo común) y \texttt{pesos\_vecino},
\texttt{pesos\_triangulos}, \texttt{pesos\_bilineal}, \texttt{pesos\_bicubica},
\texttt{pesos\_idw}, \texttt{pesos\_kriging}; \texttt{proyectar\_origen\_nacional} (EPSG:9377). \\
Kriging ordinario & \texttt{ajustar\_variograma} (semivariograma y covariograma
experimentales, ajuste por $\mathrm{SSR_w}$), \texttt{fminsearch} (Nelder--Mead equivalente al de MATLAB),
\texttt{resolver} (eliminación gaussiana con pivoteo). \\
Modelo y selección & \texttt{ajustar\_modelo} (ecuaciones~\eqref{eq:modelo18}--\eqref{eq:r2aj}),
\texttt{procesar\_malla} (seis métodos y mayor $R^2_{aj}$), \texttt{validacion\_idw} (sugerencia de $p$). \\
Flujo y salidas & \texttt{ejecutar} (flujo completo, usado por la consola y por la interfaz),
\texttt{escribir\_csv}, \texttt{escribir\_xlsx}. \\
\bottomrule
\end{tabular}
\end{table}

\begin{lstlisting}[language=Python, literate={á}{{\'a}}1 {é}{{\'e}}1 {í}{{\'i}}1 {ó}{{\'o}}1 {ú}{{\'u}}1 {ñ}{{\~n}}1 {Á}{{\'A}}1 {É}{{\'E}}1 {Í}{{\'I}}1 {Ó}{{\'O}}1 {Ú}{{\'U}}1 {Ñ}{{\~N}}1 {ü}{{\"u}}1, caption={Corrección de las alturas al sistema \emph{zero tide} (fragmento de \texttt{parametrico1.py}).}]
def delta_H_mt_a_zt(lat):
    """Ec. (1)-(2): corrección de H de mean tide a zero tide [m]."""
    s2 = math.sin(math.radians(lat)) ** 2
    return -0.198 * (1.5 * s2 - 0.5)

def delta_h_tf_a_zt(lat):
    """Ec. (3): corrección de h de tide free a zero tide [m]."""
    s2 = math.sin(math.radians(lat)) ** 2
    return 0.0602 - 0.1790 * s2 - 0.0018 * s2 * s2

def correccion_h(lat, sistema):
    """Corrección a sumar a h para llevarla a zero tide."""
    if sistema == 'zero_tide':
        return 0.0                         # ya está en zero tide: no se modifica
    if sistema == 'tide_free':
        return delta_h_tf_a_zt(lat)
    return 0.0                             # en mean tide la corteza es la misma
\end{lstlisting}

\begin{lstlisting}[language=Python, literate={á}{{\'a}}1 {é}{{\'e}}1 {í}{{\'i}}1 {ó}{{\'o}}1 {ú}{{\'u}}1 {ñ}{{\~n}}1 {Á}{{\'A}}1 {É}{{\'E}}1 {Í}{{\'I}}1 {Ó}{{\'O}}1 {Ú}{{\'U}}1 {Ñ}{{\~N}}1 {ü}{{\"u}}1, caption={Ajuste del modelo de desplazamiento e inclinación y cálculo de $R^2_{aj}$ (fragmento de \texttt{ajustar\_modelo}).}]
for p in puntos:                           # una fila de A por punto
    f, la = p['lat'], p['lon']
    filas.append([1.0, f - phi0, (la - lam0) * math.cos(math.radians(f))])
N = [[sum(r[a] * r[b] for r in filas) for b in range(u)] for a in range(u)]  # matriz normal
t = [sum(r[a] * li for r, li in zip(filas, l)) for a in range(u)]          # término independiente
x, _ = resolver(N, t)                      # parámetros: dH, x1, x2
v = [li - sum(ra * xa for ra, xa in zip(r, x)) for r, li in zip(filas, l)]  # residuos
r2 = 1 - sum(e * e for e in v) / sum((li - media) ** 2 for li in l)
r2adj = 1 - (1 - r2) * (n - 1) / (n - u)   # coeficiente de determinación ajustado
\end{lstlisting}

La interfaz (\texttt{parametrico2.py}) guía al usuario con siete pestañas, una por paso:
\begin{enumerate}
    \item guía y teoría;
    \item puntos: qué columna es la latitud, la longitud, $h$ y $H$, y en qué sistema de mareas
    vienen las alturas;
    \item mallas, con el sistema de mareas de cada una;
    \item mapa con $N$ en colores, curvas de nivel, cobertura de cada malla y ubicación de los
    puntos;
    \item parámetros y ejecución, con barra de progreso y bitácora;
    \item resultados: parámetros $\pm\sigma$, $\hat\sigma_0$, $R^2$ y $R^2_{aj}$ de cada método,
    y los puntos apartados;
    \item kriging ordinario: resumen de todos los semivariogramas locales (mediana y franja
    P10--P90) o semivariograma y covariograma de un punto.
\end{enumerate}
El cálculo corre en un hilo aparte para que la ventana no se bloquee. Los resultados se guardan
en la carpeta \texttt{resultados}:
\begin{itemize}
    \item \texttt{coordenadas\_zero\_tide.xlsx}, \texttt{puntos\_rechazados.xlsx} y
    \texttt{resumen\_modelos.xlsx};
    \item \texttt{interpolacion\_<malla>.csv}, con $N$, $l$ y $v$ de los seis métodos en cada
    punto;
    \item \texttt{reporte.txt}, con la bitácora completa.
\end{itemize}

\subsubsection{Datos y cobertura}
La Figura~\ref{fig:mallapython} muestra la malla EGM2008 de fondo ($571\times511$ nodos, paso
de $2'$), la cobertura de QGeoidCOL2023 (rectángulo punteado) y los puntos de la Red Nacional
de Nivelación. Los puntos siguen los circuitos de nivelación a lo largo de la región andina y
del norte del país, entre $\varphi=1{,}25^\circ$ y $11{,}54^\circ$ y entre
$\lambda=-77{,}37^\circ$ y $-72{,}13^\circ$. Todos quedan dentro de la cobertura de las dos
mallas, por lo que ningún punto se apartó por la regla R7. De los 
4272 registros se
aceptaron 
4210 y se apartaron 62:
\begin{itemize}
    \item 22 con altura elipsoidal nula (R4);
    \item 3 con coordenadas repetidas (R3);
    \item 2 con $|h-H|>\SI{150}{m}$ (R5);
    \item 35 por error grueso (R6). Este grupo incluye saltos de unos \SI{-22}{m} en una misma
    línea y casos con $h$ y $H$ aparentemente intercambiadas.
\end{itemize}

\begin{figure}[H]
\centering
\includegraphics[width=0.85\linewidth]{figuras/mallapython.png}
\caption{Malla EGM2008 (ondulación $N$ en colores y curvas de nivel), cobertura de
QGeoidCOL2023 (rectángulo punteado) y puntos de la Red Nacional de Nivelación (en negro).}
\label{fig:mallapython}
\end{figure}

\subsubsection{Resultados del modelo paramétrico}
Las Tablas~\ref{tab:param_qgeoid} y~\ref{tab:param_egm} presentan los parámetros del
modelo~\eqref{eq:modelo18} para cada método de interpolación, con las alturas y las mallas en
\emph{zero tide} y $n=4210$ puntos. En la bicúbica los 4210 puntos son aplicables, porque
ninguno cae en una celda de borde. El IDW se evaluó con $p=6$ y $k=16$.

\begin{table}[H]
\centering
\caption{Modelo de desplazamiento e inclinación con QGeoidCOL2023 ($n=4210$). $x_1$ y $x_2$
en m por grado. En negrita, el método de mayor $R^2_{aj}$.}
\label{tab:param_qgeoid}
\small
\begin{adjustbox}{max width=\linewidth}
\begin{tabular}{@{}lcccccc@{}}
\toprule
Método & $\delta H$ (m) & $x_1$ & $x_2$ & $\hat\sigma_0$ (m) & $R^2$ & $R^2_{aj}$ \\
\midrule
Vecino más cercano    & $0{,}1809\pm0{,}0061$ & $-0{,}0452\pm0{,}0030$ & $-0{,}0176\pm0{,}0063$ & $0{,}3939$ & $0{,}1027$ & $0{,}1023$ \\
Lineal por triángulos & $0{,}1755\pm0{,}0036$ & $-0{,}0382\pm0{,}0018$ & $-0{,}0178\pm0{,}0038$ & $0{,}2360$ & $0{,}1927$ & $0{,}1923$ \\
Bilineal              & $0{,}1731\pm0{,}0036$ & $-0{,}0375\pm0{,}0018$ & $-0{,}0170\pm0{,}0038$ & $0{,}2365$ & $0{,}1852$ & $0{,}1848$ \\
\textbf{Bicúbica}     & $\mathbf{0{,}1827\pm0{,}0036}$ & $-0{,}0396\pm0{,}0018$ & $-0{,}0169\pm0{,}0037$ & $\mathbf{0{,}2320}$ & $0{,}2056$ & $\mathbf{0{,}2053}$ \\
IDW ($p=6$)           & $0{,}1799\pm0{,}0049$ & $-0{,}0424\pm0{,}0024$ & $-0{,}0183\pm0{,}0051$ & $0{,}3100$ & $0{,}1368$ & $0{,}1364$ \\
Kriging ordinario     & $0{,}1685\pm0{,}0037$ & $-0{,}0376\pm0{,}0018$ & $-0{,}0050\pm0{,}0039$ & $0{,}2400$ & $0{,}1542$ & $0{,}1538$ \\
\bottomrule
\end{tabular}
\end{adjustbox}
\end{table}

\begin{table}[H]
\centering
\caption{Modelo de desplazamiento e inclinación con EGM2008 ($n=4210$). $x_1$ y $x_2$ en m
por grado. En negrita, el método de mayor $R^2_{aj}$.}
\label{tab:param_egm}
\small
\begin{adjustbox}{max width=\linewidth}
\begin{tabular}{@{}lcccccc@{}}
\toprule
Método & $\delta H$ (m) & $x_1$ & $x_2$ & $\hat\sigma_0$ (m) & $R^2$ & $R^2_{aj}$ \\
\midrule
Vecino más cercano    & $1{,}0715\pm0{,}0054$ & $-0{,}0931\pm0{,}0027$ & $-0{,}0264\pm0{,}0056$ & $0{,}3516$ & $0{,}3637$ & $0{,}3634$ \\
Lineal por triángulos & $1{,}0726\pm0{,}0050$ & $-0{,}0911\pm0{,}0025$ & $-0{,}0274\pm0{,}0052$ & $0{,}3258$ & $0{,}3917$ & $0{,}3914$ \\
Bilineal              & $1{,}0721\pm0{,}0050$ & $-0{,}0910\pm0{,}0025$ & $-0{,}0272\pm0{,}0052$ & $0{,}3258$ & $0{,}3909$ & $0{,}3906$ \\
\textbf{Bicúbica}     & $\mathbf{1{,}0742\pm0{,}0050}$ & $-0{,}0915\pm0{,}0025$ & $-0{,}0274\pm0{,}0052$ & $0{,}3266$ & $0{,}3926$ & $\mathbf{0{,}3923}$ \\
IDW ($p=6$)           & $1{,}0718\pm0{,}0052$ & $-0{,}0924\pm0{,}0026$ & $-0{,}0268\pm0{,}0054$ & $0{,}3375$ & $0{,}3802$ & $0{,}3799$ \\
Kriging ordinario     & $1{,}0734\pm0{,}0050$ & $-0{,}0912\pm0{,}0025$ & $-0{,}0267\pm0{,}0052$ & $0{,}3266$ & $0{,}3899$ & $0{,}3897$ \\
\bottomrule
\end{tabular}
\end{adjustbox}
\end{table}

\begin{analisisBox}{Análisis de los resultados del modelo paramétrico}
\begin{itemize}
    \item \textbf{Desplazamiento del datum.} Con EGM2008, el datum de Buenaventura queda
    \SI{1.07}{m} por encima del geoide del modelo. Este valor es del mismo orden que los
    $0{,}99\pm0{,}03$~m de \citet{guimaraes2025} para Colombia con su modelo regional y que la
    diferencia media de \SI{1.10}{m} que esos autores reportan frente a EGM2008. La diferencia de
    unos \SI{8}{cm} es razonable, porque el conjunto de puntos y el modelo de geoide no son los
    mismos. Con QGeoidCOL2023 el desplazamiento es de solo \SI{0.18}{m}, lo que indica que este
    modelo ya está referido en buena medida a la red de nivelación colombiana.
    \item \textbf{Inclinaciones.} $x_1$ es negativo con las dos mallas: la discrepancia
    disminuye hacia el norte unos \SI{4}{cm} por grado con QGeoidCOL2023 y \SI{9}{cm} por grado
    con EGM2008. Esto equivale a \SI{0.4}{m} y \SI{0.9}{m} a lo largo de los $10^\circ$ de
    latitud que cubre la red. La inclinación este--oeste ($x_2$) es menor, de \SI{-2}{cm} a
    \SI{-3}{cm} por grado. Una tendencia norte--sur de este tamaño es compatible con la
    acumulación de errores sistemáticos de la nivelación a medida que se aleja del datum, ubicado
    en el suroccidente.
    \item \textbf{Capacidad explicativa.} El $R^2_{aj}$ es bajo ($0{,}21$ con QGeoidCOL2023 y
    $0{,}39$ con EGM2008): el desplazamiento y las dos inclinaciones explican solo una parte de
    la variación de $l$. El resto es variación local, asociada sobre todo a la cordillera y a
    los errores de cada circuito. La dispersión remanente $\hat\sigma_0$ es menor con
    QGeoidCOL2023 (\SI{0.23}{m}) que con EGM2008 (\SI{0.33}{m}), lo que confirma que el modelo
    nacional describe mejor la red.
    \item \textbf{Mejor método.} La \textbf{bicúbica} obtiene el mayor $R^2_{aj}$ con las dos
    mallas, lo que coincide con la validación cruzada de la Sección~\ref{sec:cvmetodo}. Con
    QGeoidCOL2023 (paso de $5'$) el método sí importa:
    \begin{itemize}
        \item el vecino más cercano duplica la dispersión ($\hat\sigma_0=\SI{0.39}{m}$);
        \item el IDW con $p=6$, que se aproxima al vecino más cercano, también queda rezagado
        ($\hat\sigma_0=\SI{0.31}{m}$). El exponente que mejor reconstruye la propia malla no es
        necesariamente el que mejor representa $N$ en los puntos nivelados;
        \item el kriging suaviza el campo y casi anula la inclinación este--oeste.
    \end{itemize}
    Con EGM2008 (paso de $2'$) la malla es más densa y todos los métodos, salvo el vecino más
    cercano y el IDW, quedan prácticamente empatados: $R^2_{aj}$ entre $0{,}3897$ y $0{,}3923$, y
    la bilineal tiene un $\hat\sigma_0$ apenas \SI{0.8}{mm} menor que la bicúbica.
    \item \textbf{Estabilidad de $\delta H$.} El desplazamiento cambia poco con el método:
    \SI{1.4}{cm} entre extremos con QGeoidCOL2023 y \SI{0.3}{cm} con EGM2008. Su incertidumbre
    formal, de \SI{4}{mm} a \SI{6}{mm}, es pequeña por el gran número de puntos y no incluye los
    errores sistemáticos de la nivelación.
\end{itemize}
\end{analisisBox}








% =====================================================================
\section{Comparación con GeoCol2004 y EGM2008}
\label{sec:modelos}
% =====================================================================

\subsection{Magnitud y referencias de cada modelo}
Antes de restar valores de modelos distintos hay que verificar qué magnitud entrega
cada uno y respecto a qué referencia (Tabla~\ref{tab:magnitudes}).

\begin{table}[H]
\centering
\footnotesize
\setlength{\tabcolsep}{3pt}
\renewcommand{\arraystretch}{1.0}
\caption{Magnitud altimétrica, superficie, referencias, cobertura y resolución de cada modelo.}
\label{tab:magnitudes}
\begin{tabular}{@{}p{2.2cm}p{4.4cm}p{4.4cm}p{4.4cm}@{}}
\toprule
& QGeoidCOL2023 & GeoCol2004 & EGM2008 \\
\midrule
Magnitud & Anomalía de altura $\zeta$ & Ondulación geoidal $N$ & Ondulación geoidal $N$ (según la columna ``ONDU EGM08'' del archivo) \\
Superficie & Cuasigeoide & Geoide (obtenido a partir de anomalías de altura con la segunda condensación de Helmert) & Geoide respecto al elipsoide WGS84 \\
Altura asociada & Normal, $H^N=h-\zeta$ & Ortométrica, $H=h-N$ & Ortométrica, $H=h-N$ \\
Elipsoide y sistema & GRS80; IHRS ($W_0=\SI{62636853.4}{m^2/s^2}$); marea cero \citep{liu2024} & GRS80 \citep{isg}; sistema de marea no especificado en la fuente consultada; MAGNA-SIRGAS & WGS84; sistema libre de marea (\emph{tide-free}) \citep{pavlis2012}, con término de grado cero \\
Cobertura & $\ang{-4.78}$ a $\ang{12.80}$ N; $\ang{-79.30}$ a $\ang{-66.72}$ & Colombia & Global \\
Resolución & $5'$ ($\approx\SI{9}{km}$) & $2'$ ($\approx\SI{3.7}{km}$) \citep{isg} & Modelo de grado 2159 \citep{pavlis2012} ($\approx\SI{9}{km}$) \\
Desempeño declarado frente a GPS/nivelación & $\sigma=\SI{21.3}{cm}$ sin corrección; $\SI{15.8}{cm}$ con superficie de corrección de 7 parámetros (3025 puntos) \citep{liu2024} & Diferencia media de unos \SI{2}{m} con $\sigma\approx\SI{50}{cm}$ \citep{isg,sanchez2003} & $\sigma=\SI{28.1}{cm}$ con la misma corrección de 7 parámetros \citep{liu2024} \\
Obtención en P1--P5 & Esta herramienta (bilineal); contraste con MAGNA 7.0 & MAGNA 7.0 (columna ``ONDU GEOCOL2004'' del archivo) & Calculadora de UNAVCO (columna ``ONDU EGM08'' del archivo) \\
\bottomrule
\end{tabular}
\end{table}

\begin{analisisBox}{Compatibilidad conceptual}
\begin{itemize}
    \item \textbf{QGeoidCOL2023 frente a GeoCol2004}: compatibilidad \emph{parcial}. La
    diferencia $\zeta-N$ incluye la separación geoide--cuasigeoide
    (ecuación~\eqref{eq:nzeta}), además de las diferencias propias de los modelos (datos de
    gravedad, metodología, resolución y datum vertical). No debe interpretarse como ``error'' de
    ninguno de los dos modelos sin antes descontar esa separación.
    \item \textbf{QGeoidCOL2023 frente a EGM2008}: los valores del archivo son \emph{ondulaciones
    geoidales} $N$ respecto a WGS84, de modo que la compatibilidad también es \emph{parcial}. EGM2008
    permite calcular directamente la anomalía de altura (funcional \texttt{height\_anomaly} del
    servicio ICGEM, \citealp{ince2019}; \citealp{barthelmes2013}), lo que daría una comparación de
    magnitudes idénticas. Esa alternativa queda como mejora posible.
    \item \textbf{Referencias.} QGeoidCOL2023 está en el sistema de marea cero y referido al IHRS;
    las demás referencias no son necesariamente las mismas, y esa diferencia puede aportar una
    componente sistemática (Sección~\ref{sec:analisis_modelos}).
    \item Para efectos del ejercicio, QGeoidCOL2023 se adopta como referencia de
    comparación, sin considerarlo un valor verdadero.
\end{itemize}
\end{analisisBox}

\subsection{Procedimiento de obtención y verificación de los valores}
\begin{itemize}
    \item \textbf{QGeoidCOL2023.} Se interpoló $\zeta$ con la bilineal de la herramienta. El
    archivo de puntos trae también $\zeta$ redondeado al milímetro, extraído de MAGNA 7.0, y coincide con
    el valor de la herramienta a menos de \SI{0.4}{mm} (Tabla~\ref{tab:verif}).
    \item \textbf{GeoCol2004.} Se tomó la ondulación $N$ de la columna ``ONDU GEOCOL2004'' del
    archivo, extraída de MAGNA 7.0, sin recalcularla (el modelo no se incorporó a la herramienta). La altura de la columna
    ``ALT GEOCOL2004'' coincide con $h-N$ en los cinco puntos.
    \item \textbf{EGM2008.} Se tomó la ondulación $N$ de la columna ``ONDU EGM08'' del archivo, obtenida
    de la calculadora de altura geoidal de UNAVCO/EarthScope \citep{unavco} y redondeada al centímetro.
    Para verificarla de forma independiente se reprodujo esa calculadora, con el modelo \texttt{egm2008-5} (malla de $5'$ de
    GeographicLib derivada de EGM2008, \citealp{karney}) e interpolación cúbica: se ejecutó el mismo
    algoritmo (un port de \texttt{GeographicLib::Geoid}, \texttt{egm2008\_unavco.py}) y se validó contra
    los valores de referencia de GeographicLib, calculados con el modelo completo: RMS de
    \SI{4.3}{mm} en \num{20000} puntos globales y de \SI{11.3}{mm} en \num{3007} puntos en Colombia.
    Los valores del archivo difieren de esta reproducción en menos de \SI{5}{mm}
    (Tabla~\ref{tab:verif}), es decir, dentro del redondeo al centímetro, lo que confirma que la
    calculadora usa la malla \texttt{egm2008-5} con interpolación cúbica. La altura elipsoidal no
    interviene en $N$.
\end{itemize}

\begin{table}[H]
\centering
\footnotesize
\setlength{\tabcolsep}{3pt}
\renewcommand{\arraystretch}{0.9}
\caption{Verificación de los valores del archivo de puntos: $\zeta$ de QGeoidCOL2023 frente a la
interpolación bilineal de la herramienta, y $N$ de EGM2008 frente a la reproducción de la
calculadora de UNAVCO (egm2008-5, cúbica). $\Delta$ = archivo $-$ cálculo.}
\label{tab:verif}
\begin{tabular}{@{}lrrrrrr@{}}
\toprule
& \multicolumn{3}{c}{QGeoidCOL2023 $\zeta$ (m)} & \multicolumn{3}{c}{EGM2008 $N$ (m)} \\
\cmidrule(lr){2-4}\cmidrule(l){5-7}
Punto & Archivo & Herramienta & $\Delta$ (mm) & Archivo & UNAVCO & $\Delta$ (mm) \\
\midrule
P1 & \num{-13.282} & \num{-13.2822} & \num{0.2} & \num{-13.25} & \num{-13.2547} & \num{4.7} \\
P2 & \num{2.533} & \num{2.5333} & \num{-0.3} & \num{2.30} & \num{2.2999} & \num{0.1} \\
P3 & \num{11.764} & \num{11.7637} & \num{0.3} & \num{11.97} & \num{11.9717} & \num{-1.7} \\
P4 & \num{23.360} & \num{23.3603} & \num{-0.3} & \num{23.13} & \num{23.1263} & \num{3.7} \\
P5 & \num{0.826} & \num{0.8260} & \num{0.0} & \num{0.13} & \num{0.1281} & \num{1.9} \\
\bottomrule
\end{tabular}
\end{table}

\subsection{Tabla comparativa}
\begin{table}[H]
\centering
\caption{Comparación entre modelos ($\zeta$ de QGeoidCOL2023 por interpolación bilineal;
$N$ de GeoCol2004 y de EGM2008 tomados del archivo de puntos).}
\label{tab:modelos}
\small
\begin{tabular}{@{}lrrrrrr@{}}
\toprule
Punto & $\varphi$ & $\lambda$ & $h$ (m) & QGeoidCOL2023 & GeoCol2004 & EGM2008 \\
& (grados) & (grados) & & $\zeta$ (m) & $N$ (m) & $N$ (m) \\
\midrule
P1 & \num{11.45196} & \num{-72.57443} & \num{35.134} & \num{-13.282} & \num{-12.549} & \num{-13.250} \\
P2 & \num{8.43110} & \num{-75.55301} & \num{98.266} & \num{2.533} & \num{3.531} & \num{2.300} \\
P3 & \num{7.31106} & \num{-72.87202} & \num{3228.966} & \num{11.764} & \num{15.836} & \num{11.970} \\
P4 & \num{4.61497} & \num{-74.06287} & \num{2725.381} & \num{23.360} & \num{26.659} & \num{23.130} \\
P5 & \num{6.27604} & \num{-71.65145} & \num{185.792} & \num{0.826} & \num{2.410} & \num{0.130} \\
\bottomrule
\end{tabular}

\vspace{0.3cm}
\begin{tabular}{@{}lrrrrrr@{}}
\toprule
& \multicolumn{3}{c}{QGeoid $-$ GeoCol2004} & \multicolumn{3}{c}{QGeoid $-$ EGM2008} \\
\cmidrule(lr){2-4}\cmidrule(l){5-7}
Punto & m & cm & mm & m & cm & mm \\
\midrule
P1 & \num{-0.733} & \num{-73.3} & \num{-733.2} & \num{-0.032} & \num{-3.2} & \num{-32.2} \\
P2 & \num{-0.998} & \num{-99.8} & \num{-997.7} & \num{0.233} & \num{23.3} & \num{233.3} \\
P3 & \num{-4.072} & \num{-407.2} & \num{-4072.3} & \num{-0.206} & \num{-20.6} & \num{-206.3} \\
P4 & \num{-3.299} & \num{-329.9} & \num{-3298.7} & \num{0.230} & \num{23.0} & \num{230.3} \\
P5 & \num{-1.584} & \num{-158.4} & \num{-1584.0} & \num{0.696} & \num{69.6} & \num{696.0} \\
\midrule
Media & \num{-2.137} & \num{-213.7} & \num{-2137.2} & \num{0.184} & \num{18.4} & \num{184.2} \\
RMS & \num{2.510} & \num{251.0} & \num{2510.3} & \num{0.357} & \num{35.7} & \num{356.5} \\
\bottomrule
\end{tabular}
\end{table}

\begin{figure}[htbp]
\centering
\begin{tikzpicture}
\begin{axis}[
    width=0.95\linewidth, height=7.5cm,
    ybar,
    bar width=14pt,
    xmin=0.4, xmax=5.6,
    ymin=-500, ymax=100,
    xtick={1,2,3,4,5},
    xticklabels={P1,P2,P3,P4,P5},
    ytick={-500,-400,-300,-200,-100,0,100},
    ymajorgrids=true,
    grid style={gray!25},
    extra y ticks={0},
    extra y tick labels={},
    extra y tick style={grid=major, grid style={black!70, thin, solid}},
    ylabel={Diferencia QGeoidCOL2023 $-$ modelo (cm)},
    nodes near coords={\pgfmathprintnumber[fixed,precision=0]{\pgfplotspointmeta}},
    nodes near coords style={font=\scriptsize},
    legend style={at={(0.5,-0.12)}, anchor=north, legend columns=-1, draw=none, /tikz/column sep=8pt},
    legend image code/.code={\draw[#1] (0cm,-0.1cm) rectangle (0.3cm,0.1cm);},
    axis lines=left,
]
\addplot[fill=blue!60!black, draw=none] coordinates
    {(1,-73.3) (2,-99.8) (3,-407.2) (4,-329.9) (5,-158.4)};
\addplot[fill=orange, draw=none] coordinates
    {(1,-3.2) (2,23.3) (3,-20.6) (4,23.0) (5,69.6)};
\legend{QGeoid $-$ GeoCol2004, QGeoid $-$ EGM2008}
\end{axis}
\end{tikzpicture}
\caption{Diferencia entre la anomalía de altura $\zeta$ de QGeoidCOL2023 (por
interpolación bilineal) y la ondulación $N$ de GeoCol2004 y EGM2008 en los puntos
P1--P5. Medias: GeoCol2004 $-213{,}7$~cm, EGM2008 $18{,}4$~cm. RMS: GeoCol2004
$251{,}0$~cm, EGM2008 $35{,}7$~cm. Las barras más cortas indican mayor acuerdo
con QGeoidCOL2023.}
\label{fig:tabla17}
\end{figure}


\subsection{Efecto sobre la altura y control con la nivelación}
\label{sec:control}
Si en un punto se usa otro modelo en lugar de QGeoidCOL2023, la altura resultante cambia en
$H_{\text{otro}}-H^N=(h-N_{\text{otro}})-(h-\zeta)=\zeta-N_{\text{otro}}$, es decir, en la misma
magnitud de la diferencia entre modelos. Además, como el archivo trae la altura $H$ de la
nivelación geométrica, la Tabla~\ref{tab:efecto_altura} compara la altura que resulta de cada
modelo con ella. La diferencia $H^N-H=(h-\zeta)-H$ es la misma diferencia GPS/nivelación $(h-H)-\zeta$ con que
\citet{liu2024} validan el modelo, y mezcla el error del modelo con los de la altura elipsoidal, los de la nivelación, el tipo de altura de la
nivelación (nivelada y ajustada, ni ortométrica ni normal) y la inconsistencia de datum entre el IHRS y el datum
vertical local.

\begin{table}[H]
\centering
\caption{Altura nivelada $H$ y alturas obtenidas con cada modelo: $H^N=h-\zeta$ (QGeoidCOL2023),
$H_{\text{GC}}=h-N_{\text{GeoCol2004}}$ y $H_{\text{EGM}}=h-N_{\text{EGM2008}}$; diferencias
frente a la nivelada en cm.}
\label{tab:efecto_altura}
\small
\begin{tabular}{@{}lrrrrrrr@{}}
\toprule
& $H$ nivelada & $H^N$ & $H_{\text{GC}}$ & $H_{\text{EGM}}$ & $H^N-H$ & $H_{\text{GC}}-H$ & $H_{\text{EGM}}-H$ \\
Punto & (m) & (m) & (m) & (m) & (cm) & (cm) & (cm) \\
\midrule
P1 & \num{48.434} & \num{48.416} & \num{47.683} & \num{48.384} & \num{-1.8} & \num{-75.1} & \num{-5.0} \\
P2 & \num{95.573} & \num{95.733} & \num{94.735} & \num{95.966} & \num{15.9} & \num{-83.8} & \num{39.3} \\
P3 & \num{3216.550} & \num{3217.202} & \num{3213.130} & \num{3216.996} & \num{65.2} & \num{-342.0} & \num{44.6} \\
P4 & \num{2702.640} & \num{2702.021} & \num{2698.722} & \num{2702.251} & \num{-61.9} & \num{-391.8} & \num{-38.9} \\
P5 & \num{184.956} & \num{184.966} & \num{183.382} & \num{185.662} & \num{1.0} & \num{-157.4} & \num{70.6} \\
\midrule
Media (cinco puntos) & & & & & \num{3.7} & \num{-210.0} & \num{22.1} \\
RMS (cinco puntos) & & & & & \num{40.9} & \num{248.2} & \num{44.8} \\
RMS (P1, P2, P5) & & & & & \num{9.3} & \num{111.7} & \num{46.7} \\
\bottomrule
\end{tabular}
\end{table}

\begin{figure}[H]
\centering
\includegraphics[width=\linewidth]{figuras/control_nivelacion.png}
\caption{Altura de cada modelo menos la altura nivelada (cm). A la izquierda, los tres modelos;
a la derecha, QGeoidCOL2023 y EGM2008 con una escala ampliada.}
\label{fig:control}
\end{figure}

\begin{figure}[H]
\centering
\begin{subfigure}[t]{0.32\linewidth}\centering
\includegraphics[width=\linewidth]{figuras/dif_QGeoid_GeoCol2004.png}
\caption{QGeoid $-$ GeoCol2004 (cm).}
\end{subfigure}
\hfill
\begin{subfigure}[t]{0.32\linewidth}\centering
\includegraphics[width=\linewidth]{figuras/dif_QGeoid_EGM2008.png}
\caption{QGeoid $-$ EGM2008 (cm).}
\end{subfigure}
\hfill
\begin{subfigure}[t]{0.32\linewidth}\centering
\includegraphics[width=\linewidth]{figuras/dif_modelos_barras.png}
\caption{Diferencias por punto (cm).}
\end{subfigure}
\caption{Comportamiento espacial de las diferencias entre modelos sobre las isolíneas de
$\zeta$ (figuras exportadas desde la aplicación). Cada panel usa su propia escala de color.}
\label{fig:modelos}
\end{figure}

% =====================================================================
\section{Análisis}
\label{sec:analisis}
% =====================================================================

\subsection{Diferencias entre métodos}
\begin{itemize}
    \item \textbf{El vecino más cercano no es aceptable para nivelación GNSS.} Se aparta de
    la bilineal hasta \SI{25}{cm} (P5) y su RMSE en la validación cruzada es de
    \SI{48}{cm}, con errores de hasta \SI{2.5}{m}. Su error depende de la distancia
    al nodo y del gradiente local, no de la calidad del modelo.
    \item \textbf{Bilineal y triángulos son prácticamente equivalentes} (RMSE de \SI{86.6}{mm}
    y \SI{88.7}{mm}; diferencias de \SIrange{0.6}{23}{mm} en los puntos). Ambos son lineales
    por tramos y difieren solo en cómo reparten el peso dentro de la celda.
    \item \textbf{La bicúbica es la más exacta} (RMSE \SI{58}{mm}, un 33\,\% menos que la
    bilineal), porque un polinomio cúbico local sigue la curvatura del campo, que en zonas
    de montaña es importante. Su costo es que no aplica en la celda de borde y que usa pesos
    negativos, que pueden producir sobreoscilaciones.
    \item \textbf{IDW es el peor de los métodos continuos} (RMSE \SI{212}{mm}). Siempre
    promedia: el valor queda dentro del rango de los nodos vecinos y el campo se ``aplana''
    hacia la media local, lo que es un sesgo sistemático cuando $\zeta$ tiene gradiente. Sus
    diferencias frente a la bilineal fueron negativas en los cinco puntos (de \SI{-16.8}{mm} a
    \SI{-35.6}{mm}).
    \item \textbf{El kriging, con la configuración por defecto, no supera a la bilineal}
    (RMSE \SI{136}{mm}). Con una ventana más pequeña la iguala
    (Sección~\ref{sec:sensibilidad}), pero a un costo de cómputo mucho mayor. El variograma local suele resultar gaussiano con pepita casi nula, un
    modelo muy suave que con 16 vecinos produce pesos negativos y sistemas mal condicionados
    (\texttt{rcond} de $10^{-10}$ a $10^{-9}$ en P1, P2 y P5). Además, su $\sigma_{KO}$ no está calibrada: en
    P1, P2 y P5 da \SIrange{0.4}{0.7}{mm}, mientras que en P3 y P4 da unos \SI{0.28}{m}.
    Ninguno de los dos valores es coherente con el RMSE observado, así que no debe tomarse
    como la incertidumbre de $\zeta$.
\end{itemize}

\subsection{Comportamiento espacial y efecto de borde}
Las diferencias entre métodos crecen con el gradiente y la curvatura locales de $\zeta$. Sin
contar el vecino más cercano, la dispersión entre métodos es de \SIrange{21}{31}{mm} en P5 (Orinoquía),
P2 (Córdoba) y P1 (Guajira) y mayor en la cordillera Oriental (P3, \SI{83}{mm}; P4, \SI{48}{mm}).
Con un paso de $\approx\SI{9}{km}$, la malla no resuelve completamente las variaciones de
$\zeta$ causadas por la topografía abrupta. En el borde de la malla, la bicúbica deja de
ser aplicable, y el IDW y el kriging pasan a usar nodos de un solo lado, lo que aumenta
sus errores. Ninguno de los cinco puntos de control está cerca del borde, por lo que este efecto no
se observa en ellos, y la herramienta no extrapola en ningún caso. La
muestra, además, no incluye el sur del país (Amazonía, Pacífico, Nariño): en esas zonas el
comportamiento de los métodos y de los modelos podría ser distinto.

\subsection{Sensibilidad de parámetros}
La Tabla~\ref{tab:sensibilidad} muestra que IDW y kriging dependen mucho de sus
parámetros, y que los valores por defecto ($p=2$, $k=16$; ventana de 6 nodos) no son los
mejores para este campo:
\begin{itemize}
    \item En \textbf{IDW}, el error baja al aumentar $p$ y al reducir $k$ (de \SI{530}{mm}
    con $p=1$ y $k=32$ a \SI{92}{mm} con $p=4$ y $k=4$). Ambos cambios hacen el método más
    local: con $p$ alto y pocos nodos, el IDW se comporta casi como una interpolación entre
    las esquinas de la celda. Ni siquiera la mejor configuración alcanza a la bilineal, por
    el sesgo de promediado que tiene todo IDW.
    \item En \textbf{kriging}, el factor dominante es el tamaño de la ventana del
    variograma: con semiancho de 4 nodos ($\approx\pm\SI{37}{km}$) el RMSE baja a
    \SI{85.7}{mm}, al nivel de la bilineal y con un MAE menor (\SI{38}{mm} frente a
    \SI{48}{mm}). Con ventanas grandes, el variograma incorpora la tendencia regional de
    $\zeta$ y deja de describir la variabilidad local. El número de vecinos tiene un efecto
    menor, con un óptimo cerca de $n=16$.
    \item Estos óptimos se obtuvieron con un paso de prueba de $10'$. Con el paso real de
    $5'$ la escala relevante es la mitad, así que sirven como indicación de la tendencia y no
    como valores definitivos. Por la misma razón, se mantuvieron los valores por defecto para
    los resultados principales y la sensibilidad se reporta por separado.
\end{itemize}

\subsection{Diferencias entre QGeoidCOL2023 y GeoCol2004}
\label{sec:analisis_modelos}
La diferencia $\zeta-N$ (Tabla~\ref{tab:modelos}) es negativa en los cinco puntos y va de
\SI{-0.73}{m} a \SI{-4.07}{m}, con media de \SI{-2.14}{m} y RMS de \SI{2.51}{m}: GeoCol2004 entrega
ondulaciones mayores y, por tanto, alturas menores que QGeoidCOL2023 (Tabla~\ref{tab:efecto_altura}). Sus
características son las siguientes:
\begin{itemize}
    \item \textbf{No se explica por la separación geoide--cuasigeoide.} En los puntos bajos (P1, P2 y
    P5, con $H<\SI{200}{m}$) la ecuación~\eqref{eq:nzeta} da separaciones de pocos centímetros como
    máximo (del orden de \SIrange{1}{2}{cm} por cada \SI{100}{mGal} de anomalía de Bouguer, y
    se esperan anomalías menores de unos \SI{100}{mGal} en valor absoluto en esas zonas bajas),
    mientras que las diferencias son de \SIrange{0.73}{1.58}{m}. En P3 y P4 ($H\approx\SI{3200}{m}$ y
    $\SI{2700}{m}$), con una anomalía de Bouguer de $-\SI{150}{mGal}$, valor típico de la cordillera
    Oriental que se supone aquí y no se calculó, la ecuación da $\zeta-N\approx\SI{+0.4}{m}$, de signo
    contrario a las diferencias observadas (\SI{-4.07}{m} y \SI{-3.30}{m}).
    \item \textbf{Es coherente con un sesgo propio de GeoCol2004.} Frente a la nivelación, GeoCol2004 queda entre
    \SI{0.75}{m} y \SI{3.92}{m} por debajo, con una media de \SI{-2.10}{m}
    (Tabla~\ref{tab:efecto_altura}), del orden de la diferencia media de unos \SI{2}{m} que se
    documenta para ese modelo frente a puntos GPS/nivelación \citep{isg,sanchez2003}. En cambio, QGeoidCOL2023 coincide con la
    nivelación a \SI{2}{cm} o menos en P1 y P5. Con solo cinco puntos esto es un indicio y no una
    estimación del sesgo del modelo.
    \item \textbf{GeoCol2004 también se aparta de EGM2008}, en el mismo sentido y por magnitudes
    comparables: $N_{\text{GC}}-N_{\text{EGM}}$ es de \SI{+0.70}{m}, \SI{+1.23}{m}, \SI{+3.87}{m},
    \SI{+3.53}{m} y \SI{+2.28}{m} en P1 a P5. Esas diferencias crecen con la altura elipsoidal en los
    cinco puntos (de \SI{35}{m} a \SI{3229}{m}), lo que sugiere un error de GeoCol2004 que depende de la
    topografía. Es una hipótesis, que habría que confirmar con más puntos; es coherente con que
    GeoCol2004 se base en un modelo global de grado 200, un modelo digital de terreno de
    $\approx\SI{1}{km}$ y una malla de $2'$ \citep{isg}.
\end{itemize}
\noindent\textbf{Efecto sobre la altura.} Como $H_{\text{GeoCol2004}}-H^N=\zeta-N$,
cambiar de modelo altera la altura del punto en esa misma magnitud: de \SI{0.7}{m} en el Caribe a
\SI{4}{m} en la cordillera Oriental. Una altura obtenida con GeoCol2004 y una obtenida con
QGeoidCOL2023 no son intercambiables, y mezclarlas en un mismo proyecto introduce errores
métricos.

\subsection{Diferencias entre QGeoidCOL2023 y EGM2008}
La diferencia $\zeta-N_{\text{EGM2008}}$ va de \SI{-20.6}{cm} a \SI{+69.6}{cm}, con media de
\SI{+18.4}{cm}, desviación estándar de \SI{34.1}{cm} y RMS de \SI{35.7}{cm}
(Tabla~\ref{tab:modelos}):
\begin{itemize}
    \item \textbf{P1 (Guajira): \SI{-3.2}{cm}.} A baja altitud y cerca del mar los dos modelos
    coinciden a pocos centímetros.
    \item \textbf{P3 y P4 (cordillera Oriental): \SI{-20.6}{cm} y \SI{+23.0}{cm}.} Tienen signos
    opuestos y una magnitud cerca de un orden menor que las diferencias con GeoCol2004
    (\SI{-4.07}{m} y \SI{-3.30}{m}). Para P4, la
    separación geoide--cuasigeoide supuesta arriba ($\zeta-N\approx\SI{+0.4}{m}$) iría en el mismo
    sentido, pero no para P3.
    \item \textbf{P2 y P5 (\SI{+23.3}{cm} y \SI{+69.6}{cm}).} A baja altitud la separación
    $N-\zeta$ es despreciable, así que estas diferencias son de modelo o de referencia. Frente a la
    nivelación (Tabla~\ref{tab:efecto_altura}), en esos puntos la altura con QGeoidCOL2023 queda a
    \SI{15.9}{cm} y \SI{1.0}{cm} y la de EGM2008 a \SI{39.3}{cm} y \SI{70.6}{cm}. Por tanto, estas diferencias
    se deben sobre todo a EGM2008 en esos sitios, lo cual es coherente con la desviación estándar de
    \SI{28.1}{cm} que \citet{liu2024} reportan para EGM2008 frente a \SI{15.8}{cm} para QGeoidCOL2023, con la
    misma corrección y los mismos puntos GPS/nivelación.
    \item \textbf{Referencias.} QGeoidCOL2023 está en el sistema de marea cero y referido al IHRS
    ($W_0$ convencional), mientras que EGM2008 está en el sistema libre de marea y sus ondulaciones
    respecto a WGS84 incluyen un término de grado cero de $\approx\SI{-41}{cm}$, según la documentación
    del modelo. Estas diferencias de convención pueden aportar una componente sistemática, de centímetros
    a decímetros, a las diferencias entre los dos modelos. En este trabajo no se cuantificó ni se corrigió, y
    se declara como una limitación.
\end{itemize}
\noindent\textbf{Efecto sobre la altura.} Usar EGM2008 en lugar de QGeoidCOL2023 cambia la
altura entre \SI{-21}{cm} y \SI{+70}{cm} (Tabla~\ref{tab:modelos}). El error de
interpolación de la malla egm2008-5 (RMS de \SI{1.1}{cm} en Colombia) es despreciable frente
a esas diferencias, que provienen de los modelos y no del cálculo.

\subsection{Coherencia altimétrica con la nivelación}
\label{sec:analisis_control}
La Tabla~\ref{tab:efecto_altura} y la Figura~\ref{fig:control} confrontan las alturas de cada modelo con la de
la nivelación geométrica:
\begin{itemize}
    \item \textbf{Puntos bajos (P1, P2 y P5; $H<\SI{200}{m}$).} Aquí, sea ortométrica o normal la altura de la nivelación (el archivo solo
    indica que es nivelada y ajustada), la diferencia entre ambos tipos es de pocos centímetros, así que la
    comparación es directa. QGeoidCOL2023 queda a
    \SI{-1.8}{cm}, \SI{+15.9}{cm} y \SI{+1.0}{cm} (RMS de \SI{9.3}{cm}), EGM2008 a
    \SI{-5.0}{cm}, \SI{+39.3}{cm} y \SI{+70.6}{cm} (\SI{46.7}{cm}) y GeoCol2004 a
    \SI{-75.1}{cm}, \SI{-83.8}{cm} y \SI{-157.4}{cm} (\SI{111.7}{cm}).
    \item \textbf{Puntos de montaña (P3 y P4).} Las diferencias de QGeoidCOL2023 son de \SI{+65.2}{cm} y
    \SI{-61.9}{cm}, y las de EGM2008 de \SI{+44.6}{cm} y \SI{-38.9}{cm}: del mismo signo y de magnitud
    parecida, lo que apunta a un componente común a ambos modelos de alta resolución, o a los datos de
    esos puntos (nivelaciones de 2005 y 2001; alturas elipsoidales GNSS), más que a un defecto
    específico de QGeoidCOL2023. Si la altura nivelada fuera ortométrica, la separación $H^N-H=N-\zeta$ (del orden de
    $-\SI{0.4}{m}$ con la anomalía de Bouguer supuesta) podría explicar parte de P4 pero no de P3; como el
    archivo no establece el tipo de altura, esa explicación no se usa aquí.
    \item \textbf{Contexto.} En la validación de \citet{liu2024}, las diferencias entre
    QGeoidCOL2023 y los 3025 puntos GPS/nivelación tienen media de \SI{+19.2}{cm}, desviación estándar de
    \SI{21.3}{cm} y un rango de \SI{-46.1}{cm} a \SI{+87.7}{cm} antes de la corrección. Las diferencias de los cinco
    puntos (de \SI{-61.9}{cm} a \SI{+65.2}{cm}, media de \SI{+3.7}{cm}) son del mismo orden, salvo P4, que
    queda \SI{16}{cm} por debajo del mínimo reportado. Con cinco puntos no es posible estimar la exactitud
    del modelo, y no se sabe si estos puntos forman parte del conjunto de validación de ese trabajo, por lo que
    no se consideran independientes.
\end{itemize}
\noindent Estos resultados no establecen la superioridad absoluta de ningún modelo: con cinco puntos, la
nivelación y el GNSS tienen sus propios errores, y el datum vertical local no coincide con el IHRS. Sí
muestran que, en este conjunto, QGeoidCOL2023 reproduce la nivelación mejor que EGM2008 y mucho mejor
que GeoCol2004, y que la elección del modelo cambia la altura entre decímetros y metros, mientras
que la del método de interpolación la cambia como máximo unos \SI{8}{cm} (sin el vecino más cercano).

% =====================================================================
\section{Conclusiones}
% =====================================================================
\begin{itemize}
    \item La malla QGeoidCOL2023 se validó con controles explícitos y documentados: es una
    malla completa, rectangular y regular de $212\times152$ nodos con paso de $5'$, sin
    datos inválidos, duplicados ni faltantes.
    \item Los seis métodos se implementaron con un núcleo común ($\hat\zeta=\sum w_i\zeta_i$)
    que expone nodos, pesos y contribuciones y verifica $\sum w_i=1$ en cada cálculo. Las
    implementaciones bilineal, bicúbica y de vecino más cercano coinciden con \texttt{interp2}
    a nivel de $10^{-14}$\,m. Ningún método extrapola, y la bicúbica declara su falta de
    soporte en el borde en lugar de sustituirse.
    \item En los puntos de control, la elección del método cambia $\zeta$ (y por tanto
    $H^N$) entre \SI{21}{mm} y \SI{83}{mm} entre los métodos continuos, y hasta \SI{25}{cm} con
    el vecino más cercano. Las diferencias son mayores donde el campo tiene más gradiente
    y curvatura.
    \item La validación cruzada por decimación ordena los métodos así: bicúbica, bilineal
    $\approx$ triángulos, kriging, IDW y vecino más cercano. La bilineal sigue siendo una
    elección razonable por su simplicidad y trazabilidad, y la bicúbica es preferible donde
    tiene soporte.
    \item La diferencia $\zeta-N$ entre QGeoidCOL2023 y GeoCol2004 es negativa en los cinco puntos
    (de \SI{-0.7}{m} a \SI{-4.1}{m}, media de \SI{-2.1}{m}). Esa diferencia no se explica por la separación geoide--cuasigeoide sino por el
    sesgo de GeoCol2004, que frente a la nivelación es del orden de lo documentado para ese modelo (unos
    \SI{2}{m}). La diferencia frente a la nivelación de QGeoidCOL2023 es de \SI{2}{cm} o menos en P1 y P5.
    \item Frente a EGM2008 las diferencias son de \SI{-21}{cm} a \SI{+70}{cm} (RMS de \SI{36}{cm}). En los
    puntos bajos donde EGM2008 se aparta más (P2 y P5) la altura con QGeoidCOL2023 es la que más se acerca a
    la nivelación, de modo que la diferencia se atribuye principalmente a EGM2008, con la reserva de las
    posibles diferencias de referencia (sistema de marea y término de grado cero), que no se cuantificaron.
    \item En los dos puntos de montaña, QGeoidCOL2023 y EGM2008 difieren de la nivelación en decímetros y
    con el mismo signo. Esto indica un componente común a ambos modelos o a los datos de esos puntos, y no
    permite concluir sobre la exactitud de QGeoidCOL2023 en zonas montañosas.
    \item La elección del modelo influye sobre la altura mucho más (decímetros a metros) que la del
    método de interpolación (centímetros). En consecuencia, $N$ y $\zeta$ (y $H$ y $H^N$) no son
    intercambiables.
    \item Limitaciones: cinco puntos (todos al norte de \ang{4.6} N, sin puntos en la Amazonía, el Pacífico ni el
    sur) no bastan para estimar la exactitud de los modelos; los valores de GeoCol2004 y EGM2008 provienen del
    archivo de puntos (MAGNA 7.0 y la calculadora de UNAVCO; EGM2008 se verificó de forma independiente); y las
    alturas niveladas (ajustadas, de tipo no especificado) no se reducen a un mismo datum vertical con el IHRS.
\end{itemize}

% =====================================================================
\begin{thebibliography}{9}
% =====================================================================
\bibitem[Olarte(2025)]{olarte2025}
Olarte Ramírez, H. (2025). \textit{Estimación de un modelo micro-cuasigeoidal mediante el
modelamiento del campo de gravedad local para la incorporación de la primera estación de
Colombia en el IHRF}. Trabajo de grado para optar al título de Ingeniero Catastral y Geodesta.
Director: C.~E. Melo Martínez, Ph.D. Facultad de Ingeniería, Ingeniería Catastral y Geodesia,
Universidad Distrital Francisco José de Caldas, Bogotá, Colombia.

\bibitem[Barthelmes(2013)]{barthelmes2013}
Barthelmes, F. (2013). \textit{Definition of functionals of the geopotential and their
calculation from spherical harmonic models}. Scientific Technical Report STR09/02
(revisado 2013). GFZ German Research Centre for Geosciences, Potsdam.
\url{https://doi.org/10.2312/GFZ.b103-0902-26}

\bibitem[Cressie(1993)]{cressie1993}
Cressie, N. (1993). \textit{Statistics for Spatial Data} (ed.\ revisada). John Wiley \& Sons.

\bibitem[Hofmann-Wellenhof y Moritz(2006)]{hofmannmoritz2006}
Hofmann-Wellenhof, B., Moritz, H. (2006). \textit{Physical Geodesy} (2.ª ed.). Springer.

\bibitem[Ince et~al.(2019)]{ince2019}
Ince, E.S., Barthelmes, F., Reißland, S., Elger, K., Förste, C., Flechtner, F., Schuh, H.
(2019). ICGEM -- 15 years of successful collection and distribution of global gravitational
models, associated services, and future plans. \textit{Earth System Science Data}, 11,
647--674. \url{https://doi.org/10.5194/essd-11-647-2019}

\bibitem[Keys(1981)]{keys1981}
Keys, R. (1981). Cubic convolution interpolation for digital image processing.
\textit{IEEE Transactions on Acoustics, Speech, and Signal Processing}, 29(6), 1153--1160.

\bibitem[Pavlis et~al.(2012)]{pavlis2012}
Pavlis, N.K., Holmes, S.A., Kenyon, S.C., Factor, J.K. (2012). The development and
evaluation of the Earth Gravitational Model 2008 (EGM2008). \textit{Journal of Geophysical
Research}, 117, B04406. \url{https://doi.org/10.1029/2011JB008916}

\bibitem[Shepard(1968)]{shepard1968}
Shepard, D. (1968). A two-dimensional interpolation function for irregularly-spaced
data. \textit{Proceedings of the 1968 ACM National Conference}, 517--524.

\bibitem[Snyder(1987)]{snyder1987}
Snyder, J.P. (1987). \textit{Map Projections -- A Working Manual}. USGS Professional
Paper 1395.

\bibitem[Torge y Müller(2012)]{torgemuller2012}
Torge, W., Müller, J. (2012). \textit{Geodesy} (4.ª ed.). De Gruyter.

\bibitem[Karney(s.f.)]{karney}
Karney, C.F.F. (s.f.). \textit{GeographicLib: Geoid height}. Documentación de la clase
\texttt{GeographicLib::Geoid} y de las mallas \texttt{egm2008-5}.
\url{https://geographiclib.sourceforge.io/html/geoid.html}

\bibitem[UNAVCO/EarthScope(s.f.)]{unavco}
UNAVCO/EarthScope Consortium (s.f.). \textit{Geoid Height Calculator}.
\url{https://www.unavco.org/software/geodetic-utilities/geoid-height-calculator/geoid-height-calculator.html}
Consultado el 6 de octubre de 2026.

\bibitem[IGAC(2026)]{igacnoticia}
Instituto Geográfico Agustín Codazzi (2026). \textit{Después de 20 años, Colombia cuenta con un
nuevo modelo oficial de referencia de alturas: QGeoidCOL2023}. Noticia institucional.
\url{https://www.igac.gov.co/noticias/despues-de-20-anos-colombia-cuenta-con-un-nuevo-modelo-oficial-de-referencia-de-alturas-qgeoidcol2023}
Consultado el 9 de octubre de 2026.

\bibitem[ISG(s.f.)]{isg}
International Service for the Geoid (s.f.). \textit{Fichas de los modelos GEOCOL2004 y QgeoidCOL2023}.
Politecnico di Milano. \url{https://isgeoid.polimi.it/Geoid/America/Colombia/colombia04_g.html} y
\url{https://isgeoid.polimi.it/Geoid/America/Colombia/colombia23_g.html}
Consultado el 9 de octubre de 2026.

\bibitem[Liu et~al.(2024)]{liu2024}
Liu, Q., Schmidt, M., Sánchez, L., Moisés, L., Cortez, D. (2024). Modelado regional del campo de
gravedad de alta resolución en regiones con desafíos de datos para la realización de sistemas de
alturas basados en geopotencial (documento técnico de QGeoidCOL2023, traducción al español publicada
por el IGAC). \textit{Earth, Planets and Space}, 76, 35.
\url{https://doi.org/10.1186/s40623-024-01981-1}

\bibitem[Sánchez(2003)]{sanchez2003}
Sánchez, L. (2003). \textit{Determinación de la superficie vertical de referencia para Colombia}.
Tesis doctoral, Technische Universität Dresden (citada en la ficha de GEOCOL2004 de ISG).

\bibitem[Guimarães et~al.(2025)]{guimaraes2025}
Guimarães, G.N., Matos, A.C.O.C., Blitzkow, D., Almeida, F.G.V. (2025). High-resolution
regional gravity field modeling to connect the local vertical systems in South America to
the IHRF. \textit{Earth, Planets and Space}, 77, 98.
\url{https://doi.org/10.1186/s40623-025-02226-5}


\end{thebibliography}

% =====================================================================
\appendix
\section{Tablas de nodos de IDW y kriging en P4}
\label{anx:tablas}
% =====================================================================

\begin{table}[H]
\centering
\caption{P4: nodos y pesos de IDW ($p=2$, $k=16$; $\hat\zeta=\SI{23.3247}{m}$) y de kriging
ordinario ($n=16$; $\hat\zeta=\SI{23.3564}{m}$). Mismos 16 nodos, ordenados por distancia.}
\label{tab:nodos_idw_ko}
\small
\begin{tabular}{@{}rrrrrrrr@{}}
\toprule
& & & & \multicolumn{2}{c}{IDW} & \multicolumn{2}{c}{Kriging} \\
\cmidrule(lr){5-6}\cmidrule(l){7-8}
\# & $\varphi$ & $\lambda$ & $\zeta_i$ (m) & $w_i$ & $w_i\zeta_i$ & $\lambda_i$ & $\lambda_i\zeta_i$ \\
\midrule
1 & \num{4.633333} & \num{-74.050000} & \num{23.332} & \num{0.630354} & \num{14.7074} & \num{0.248233} & \num{5.7918} \\
2 & \num{4.550000} & \num{-74.050000} & \num{23.729} & \num{0.072414} & \num{1.7183} & \num{0.210289} & \num{4.9899} \\
3 & \num{4.633333} & \num{-74.133333} & \num{22.937} & \num{0.059550} & \num{1.3659} & \num{0.201376} & \num{4.6190} \\
4 & \num{4.550000} & \num{-74.133333} & \num{23.388} & \num{0.034461} & \num{0.8060} & \num{0.153020} & \num{3.5788} \\
5 & \num{4.633333} & \num{-73.966667} & \num{23.572} & \num{0.032907} & \num{0.7757} & \num{0.145262} & \num{3.4241} \\
6 & \num{4.716667} & \num{-74.050000} & \num{22.958} & \num{0.030236} & \num{0.6942} & \num{0.086738} & \num{1.9913} \\
7 & \num{4.550000} & \num{-73.966667} & \num{23.635} & \num{0.023466} & \num{0.5546} & \num{0.116801} & \num{2.7606} \\
8 & \num{4.716667} & \num{-74.133333} & \num{22.630} & \num{0.020713} & \num{0.4687} & \num{0.062797} & \num{1.4211} \\
9 & \num{4.716667} & \num{-73.966667} & \num{23.305} & \num{0.016162} & \num{0.3767} & \num{-0.009123} & \num{-0.2126} \\
10 & \num{4.466667} & \num{-74.050000} & \num{23.900} & \num{0.014339} & \num{0.3427} & \num{0.013617} & \num{0.3255} \\
11 & \num{4.633333} & \num{-74.216667} & \num{22.602} & \num{0.013156} & \num{0.2973} & \num{0.025824} & \num{0.5837} \\
12 & \num{4.466667} & \num{-74.133333} & \num{23.962} & \num{0.011771} & \num{0.2821} & \num{-0.040964} & \num{-0.9816} \\
13 & \num{4.550000} & \num{-74.216667} & \num{23.012} & \num{0.011332} & \num{0.2608} & \num{-0.031283} & \num{-0.7199} \\
14 & \num{4.466667} & \num{-73.966667} & \num{23.514} & \num{0.010147} & \num{0.2386} & \num{-0.057210} & \num{-1.3452} \\
15 & \num{4.633333} & \num{-73.883333} & \num{23.531} & \num{0.009690} & \num{0.2280} & \num{-0.058619} & \num{-1.3794} \\
16 & \num{4.716667} & \num{-74.216667} & \num{22.330} & \num{0.009303} & \num{0.2077} & \num{-0.066757} & \num{-1.4907} \\
\midrule
$\sum$ & & & & \num{1.000000} & \num{23.3247} & \num{1.000000} & \num{23.3564} \\
\bottomrule
\end{tabular}
\end{table}

% =====================================================================
\section{Organización del código fuente}
\label{anx:codigo}
% =====================================================================
El código completo se entrega como archivos independientes (Anexo~\ref{anx:archivos}) y no se transcribe
en este documento. La aplicación (\texttt{appQGeoidCOL2023.m}) es un archivo único con las funciones
siguientes:
\begin{itemize}
    \item \textbf{Interfaz y modos de uso:} \texttt{appQGeoidCOL2023} (interfaz gráfica),
    \texttt{appQGeoidCOL2023('lote')} (validación, seis métodos, comparación de métodos y de modelos y validación cruzada)
    y \texttt{appQGeoidCOL2023('sensibilidad')} (análisis de sensibilidad de parámetros).
    \item \textbf{Lectura y validación:} \texttt{cargarYValidarMalla}, \texttt{autoverificarImplementacion}.
    \item \textbf{Interpolación:} \texttt{interpolarPunto} (núcleo común) y \texttt{pesosVecinoMasCercano},
    \texttt{pesosTriangulos}, \texttt{pesosBilineal}, \texttt{pesosBicubica}, \texttt{pesosIDW}, \texttt{pesosKriging}.
    \item \textbf{Evaluación:} \texttt{compararMetodos}, \texttt{validacionDecimacion}, \texttt{compararModelos}.
    \item \textbf{Gráficos:} \texttt{graficarMallaGeneral}, \texttt{graficarVistaLocal},
    \texttt{graficarComparacionMetodos}, \texttt{graficarValidacionCruzada}, \texttt{graficarDiferenciasModelos}.
\end{itemize}
Los fragmentos esenciales (núcleo común y pesos de la lineal por triángulos) se muestran en la
Sección~\ref{sec:procedimiento}.


% =====================================================================
\section{Archivos de la entrega}
\label{anx:archivos}
% =====================================================================
\begin{table}[H]
\centering
\caption{Archivos que acompañan al informe.}
\small
\renewcommand{\arraystretch}{1.2}
\begin{tabular}{@{}p{5.6cm}p{9.8cm}@{}}
\toprule
Archivo & Contenido \\
\midrule
\texttt{appQGeoidCOL2023.m} & Aplicación (interfaz, modo por lotes y de sensibilidad). \\
\texttt{QGeoidCOL2023.txt} & Malla oficial de $\zeta$ (entrada). \\
\texttt{valores\_control\_P1\_P5.csv} & Puntos, alturas niveladas y valores de GeoCol2004, EGM2008 y QGeoidCOL2023 tomados de \texttt{5 PUNTOS CALCULADOS COLOMBIA.xlsx}. \\
\texttt{generarResultadosInforme.m} & Reproduce las figuras (\texttt{informe/figuras}) y los datos (\texttt{informe/datos}). \\
\texttt{egm2008\_unavco.py} & Reproducción de la calculadora de UNAVCO (egm2008-5); requiere \texttt{egm2008-5.pgm} y \texttt{GeoidHeights.dat} de GeographicLib. \\
\bottomrule
\end{tabular}
\end{table}

\end{document}
